% =====================================================================
%  fig1_objective.tex --- Figure 1: the objective and its scope
%
%  CONCEPTUAL FIGURE.  NO DATA.  No number in this figure is measured;
%  every glyph is a schematic.  Nothing here asserts a result.
%
%  Replaces the old five-band figures/fig1_method.tex, which was too
%  dense for a first-page explanatory figure and whose bands (b)/(c)
%  gave the gradient-supervision path equal visual status.  Per the
%  build brief the force branch is REMOVED from this figure entirely;
%  gradient supervision appears only later, as an experimental
%  comparator, and is not drawn here.
%
%  Pure TikZ.  No external images, no pgfplots, no \includegraphics,
%  no arrows.meta, no pgfmathrandom (byte-reproducible across compiles).
%  Okabe--Ito colour-blind-safe palette, semantics shared with
%  make_experiment_figures.py::PALETTE:
%     blue    model / flow / learned quantities
%     green   the value term (energy values, L_value)
%     orange  the teacher (universal neural potential)
%     grey    neutral structure
%  Vermilion (= gradient supervision) and purple (= scrambled control)
%  are deliberately ABSENT: neither concept appears in this figure.
%
%  USAGE
%   (a) standalone (verification / slides):   pdflatex fig1_objective.tex
%   (b) inside the paper -- main.tex already \usepackage{tikz} and
%       \usepackage{amsmath}:
%         % =====================================================================
%  fig1_objective_float.tex --- the float wrapper for Figure 1.
%
%  Figure 1 is pure TikZ, so this wrapper \input's the picture directly
%  rather than \includegraphics-ing a PDF (the same mechanism the old
%  fig1_method.tex used).  main.tex already loads tikz and amsmath.
%
%  Include from sections/01_intro.tex with:
%      % =====================================================================
%  fig1_objective_float.tex --- the float wrapper for Figure 1.
%
%  Figure 1 is pure TikZ, so this wrapper \input's the picture directly
%  rather than \includegraphics-ing a PDF (the same mechanism the old
%  fig1_method.tex used).  main.tex already loads tikz and amsmath.
%
%  Include from sections/01_intro.tex with:
%      % =====================================================================
%  fig1_objective_float.tex --- the float wrapper for Figure 1.
%
%  Figure 1 is pure TikZ, so this wrapper \input's the picture directly
%  rather than \includegraphics-ing a PDF (the same mechanism the old
%  fig1_method.tex used).  main.tex already loads tikz and amsmath.
%
%  Include from sections/01_intro.tex with:
%      \input{figures/fig1_objective_float}
%
%  The picture is drawn at 13.90cm = 5.47in, i.e. full ICLR single-column
%  width at scale 1.0.  Do NOT wrap it in \resizebox and do NOT scale it:
%  the in-figure type is \small / \footnotesize / \scriptsize of the 10pt
%  body font and scaling would desynchronise it from the caption.
%
%  PROVENANCE.  Conceptual figure, no data.  Built per the Figure 1 brief:
%  four panels (generator / teacher values / the value term / scope), and
%  the gradient-supervision branch of the superseded fig1_method.tex is
%  removed entirely rather than demoted.  No panel asserts a result.
%
%  REVISION C13 (figure-editor pass).  The framing moves to the method: panel
%  (c) is titled "Proposed energy-value supervision" and states inside the box
%  that the partition function cancels; panel (d)'s heading is "Local
%  supervision and the path to global control" and its closing line names the
%  extension (connecting groups so they share one offset) alongside the
%  unchanged qualification that b_B - b_A and the relative mass are not
%  identified.  No mathematical content was removed.
%
%  REVISION C8.  The caption is rewritten only where panels (c) and (d)
%  changed.  (c) now names the complete objective, so the caption states
%  it.  (d) now shows two disconnected groups of ONE composition, so the
%  caption no longer says the offset is free "between compositions"; the
%  freedom is already there inside a single conditional density.  The
%  verb "pins" is gone: the term CONSTRAINS THE LOCAL RELATION between
%  log-density and energy within a group.  The disconnection claim is
%  empirical and audited -- distinct parents sharing a composition supply
%  distinct perturbation clouds that share no sampled geometry (31% of
%  training parents sit on a repeated composition and are still
%  disconnected), so the overlap graph is a set of isolated vertices.
% =====================================================================
\begin{figure}[t]
  \centering
  \providecommand{\figOneScale}{1.0}\renewcommand{\figOneScale}{1.000}%
% SCALE = 1.000, THE FIGURE's NATURAL SIZE.  Do not lower it to buy page space.
% TikZ's `scale' multiplies COORDINATES but not type: every label in this
% picture is set at an absolute size (\small / \footnotesize / \scriptsize of
% the 10pt body), so any scale < 1 shrinks the panel frames out from under
% text that keeps its size, and the text runs out of the panels.
% MEASURED 2026-08-30 (600 dpi render of the standalone picture; panel frames
% located from the rendered strokes, so the calibration is self-checking at
% 236.22 px/cm on both axes):
%     scale 0.86 -> 4085 ink pixels outside the four panel frames, in 6
%                   regions.  Panel (b)'s body line runs to x = 14.59 cm,
%                   i.e. 0.71 cm past its own frame and 0.69 cm past the
%                   picture's 13.90 cm bounding box; panel (a)'s last two lines
%                   cross into panel (b)'s frame; panel (c)'s closing identity
%                   line crosses the left frame at x = -0.17 cm.
%     scale 1.000 -> 0 ink pixels outside the panels, at every ink threshold
%                   tested (lum < 128 and lum < 200, with 5 px of stroke
%                   tolerance).  All four panels are clean.
% PAGE COST: none that changes the budget.  With 1.000 the main text ends on
% page 9 and a probe \vspace inserted before \label{endofmaintext} shows a
% further +108pt fits before it spills to page 10 (+114pt spills); at 0.86 the
% same probe breaks between +108pt and +120pt.  The fix costs at most about
% one body line of final-page headroom.
  \input{figures/fig1_objective}
  \caption{\textbf{Energy-value supervision without a partition function.}
    Ideal-objective schematic; no measured quantity appears here. The archived implementation does not realise this density construction (\Cref{sec:method:impl}).
    \textbf{(a)}~With the endpoint discrete labels clamped, reversing the
    positional ODE defines the \emph{clamped positional-flow density}
    $\log q_\theta(x \mid c)$, which need not equal the joint generator's
    conditional density. \textbf{(b)}~A universal neural potential values each
    of $K$ displaced geometries of one parent. \textbf{(c)}~The proposed
    objective is flow matching plus the within-group variance of
    $\log q_\theta(x_k \mid c) + E_k/kT$, which vanishes exactly when
    $\log q_\theta = -E_k/kT + b$ inside the group: the partition function
    cancels and the offset $b$ is free. \textbf{(d)}~Two disconnected groups of
    the \emph{same} density keep independent offsets, so the term constrains
    the local density--energy relation inside each group and leaves their
    relative mass unidentified; connecting groups so they share one offset is
    the extension toward global control. Panel detail in \Cref{app:figdetail}.}
  \label{fig:objective}
\end{figure}

%
%  The picture is drawn at 13.90cm = 5.47in, i.e. full ICLR single-column
%  width at scale 1.0.  Do NOT wrap it in \resizebox and do NOT scale it:
%  the in-figure type is \small / \footnotesize / \scriptsize of the 10pt
%  body font and scaling would desynchronise it from the caption.
%
%  PROVENANCE.  Conceptual figure, no data.  Built per the Figure 1 brief:
%  four panels (generator / teacher values / the value term / scope), and
%  the gradient-supervision branch of the superseded fig1_method.tex is
%  removed entirely rather than demoted.  No panel asserts a result.
%
%  REVISION C13 (figure-editor pass).  The framing moves to the method: panel
%  (c) is titled "Proposed energy-value supervision" and states inside the box
%  that the partition function cancels; panel (d)'s heading is "Local
%  supervision and the path to global control" and its closing line names the
%  extension (connecting groups so they share one offset) alongside the
%  unchanged qualification that b_B - b_A and the relative mass are not
%  identified.  No mathematical content was removed.
%
%  REVISION C8.  The caption is rewritten only where panels (c) and (d)
%  changed.  (c) now names the complete objective, so the caption states
%  it.  (d) now shows two disconnected groups of ONE composition, so the
%  caption no longer says the offset is free "between compositions"; the
%  freedom is already there inside a single conditional density.  The
%  verb "pins" is gone: the term CONSTRAINS THE LOCAL RELATION between
%  log-density and energy within a group.  The disconnection claim is
%  empirical and audited -- distinct parents sharing a composition supply
%  distinct perturbation clouds that share no sampled geometry (31% of
%  training parents sit on a repeated composition and are still
%  disconnected), so the overlap graph is a set of isolated vertices.
% =====================================================================
\begin{figure}[t]
  \centering
  \providecommand{\figOneScale}{1.0}\renewcommand{\figOneScale}{1.000}%
% SCALE = 1.000, THE FIGURE's NATURAL SIZE.  Do not lower it to buy page space.
% TikZ's `scale' multiplies COORDINATES but not type: every label in this
% picture is set at an absolute size (\small / \footnotesize / \scriptsize of
% the 10pt body), so any scale < 1 shrinks the panel frames out from under
% text that keeps its size, and the text runs out of the panels.
% MEASURED 2026-08-30 (600 dpi render of the standalone picture; panel frames
% located from the rendered strokes, so the calibration is self-checking at
% 236.22 px/cm on both axes):
%     scale 0.86 -> 4085 ink pixels outside the four panel frames, in 6
%                   regions.  Panel (b)'s body line runs to x = 14.59 cm,
%                   i.e. 0.71 cm past its own frame and 0.69 cm past the
%                   picture's 13.90 cm bounding box; panel (a)'s last two lines
%                   cross into panel (b)'s frame; panel (c)'s closing identity
%                   line crosses the left frame at x = -0.17 cm.
%     scale 1.000 -> 0 ink pixels outside the panels, at every ink threshold
%                   tested (lum < 128 and lum < 200, with 5 px of stroke
%                   tolerance).  All four panels are clean.
% PAGE COST: none that changes the budget.  With 1.000 the main text ends on
% page 9 and a probe \vspace inserted before \label{endofmaintext} shows a
% further +108pt fits before it spills to page 10 (+114pt spills); at 0.86 the
% same probe breaks between +108pt and +120pt.  The fix costs at most about
% one body line of final-page headroom.
  % =====================================================================
%  fig1_objective.tex --- Figure 1: the objective and its scope
%
%  CONCEPTUAL FIGURE.  NO DATA.  No number in this figure is measured;
%  every glyph is a schematic.  Nothing here asserts a result.
%
%  Replaces the old five-band figures/fig1_method.tex, which was too
%  dense for a first-page explanatory figure and whose bands (b)/(c)
%  gave the gradient-supervision path equal visual status.  Per the
%  build brief the force branch is REMOVED from this figure entirely;
%  gradient supervision appears only later, as an experimental
%  comparator, and is not drawn here.
%
%  Pure TikZ.  No external images, no pgfplots, no \includegraphics,
%  no arrows.meta, no pgfmathrandom (byte-reproducible across compiles).
%  Okabe--Ito colour-blind-safe palette, semantics shared with
%  make_experiment_figures.py::PALETTE:
%     blue    model / flow / learned quantities
%     green   the value term (energy values, L_value)
%     orange  the teacher (universal neural potential)
%     grey    neutral structure
%  Vermilion (= gradient supervision) and purple (= scrambled control)
%  are deliberately ABSENT: neither concept appears in this figure.
%
%  USAGE
%   (a) standalone (verification / slides):   pdflatex fig1_objective.tex
%   (b) inside the paper -- main.tex already \usepackage{tikz} and
%       \usepackage{amsmath}:
%         \input{figures/fig1_objective_float}
%       or, if the writer prefers to own the float:
%         \begin{figure}[t]\centering
%           \input{figures/fig1_objective}
%           \caption{...}\label{fig:objective}
%         \end{figure}
%
%  SIZE.  Natural size 13.90cm x 8.80cm == 5.47in x 3.46in.
%  ICLR is a SINGLE 5.5in text column, so this is a full-text-width
%  figure at scale 1.0.  DO NOT rescale it and DO NOT wrap it in
%  \resizebox: every glyph inside is \small / \footnotesize /
%  \scriptsize of the 10pt body font, and scaling desynchronises
%  in-figure text from the caption.
%
%  LAYOUT.  A 2x2 grid, read left-to-right, top-to-bottom:
%    (a) the generator      composition AND geometry from one flow
%    (b) the teacher        energy values for K displacements of one parent
%    (c) the objective      L_ours = L_FM + lambda_E L_value, and nothing else
%    (d) the scope          two basins of ONE composition; offset free
%  Panel widths 6.80 / 6.78 cm; no wire crosses a panel boundary, so
%  there is no routing to collide with text.
%
%  REVISION C8 (this pass).  Two changes, both confined to panels (c)/(d):
%   * (c) now displays the complete objective, L_ours = L_FM + lambda_E
%     L_value, so the reader cannot infer a hidden force or anchor branch.
%     Neither branch was ever drawn here and neither is added.
%   * (d) previously drew two groups of two DIFFERENT compositions
%     (c_A, c_B).  That was wrong twice over.  It misstated the failure
%     mode -- the unidentified offset is already present WITHIN a single
%     conditional density -- and it flirted with the claim that one
%     composition contributes exactly one group, which the composition
%     audit refutes: in the training perturbation shard
%     (processed_data/omol25_4m_processed/perturbation_train_n30000_s0.pt)
%     30,000 parent groups carry only
%     22,822 distinct (element multiset, total charge) keys, 2,204 keys
%     are shared by more than one parent (9,382 parents, 31%; max
%     multiplicity 80).  Panel (d) therefore now shows two groups A and B
%     of the SAME composition c: two disconnected supervised regions of
%     one density p_theta(.|c), each with its own free offset b_A, b_B.
%     "ordering fixed" / "pinned" wording is replaced throughout by
%     "the local relation between log-density and energy is constrained".
%  In the EVALUATION population (val_data_processed.pt; the 120 grid
%  parents and the 93 primary parents) every parent has a distinct
%  composition, multiplicity 1 -- that is an audited property of the
%  evaluation set, not of the supervision design, and this figure does
%  not assert it.
% =====================================================================

% ---- standalone bootstrap (skipped when \input into a tikz-enabled doc) ----
\ifx\tikzpicture\undefined
  \documentclass[10pt]{article}
  % 0.05cm of slack past the 13.90cm bounding box so that the panel border
  % stroke width does not report a spurious overfull \hbox in standalone mode.
  \usepackage[paperwidth=14.40cm,paperheight=9.20cm,margin=0.20cm,noheadfoot]{geometry}
  \usepackage{amsmath,amssymb,bm}
  \usepackage{times}
  \usepackage{tikz}
  \setlength{\topskip}{0pt}
  \setlength{\parindent}{0pt}
  \pagestyle{empty}
  \begin{document}%
  \def\figOneFinish{\end{document}}
\fi
\providecommand{\figOneFinish}{\relax}
\providecommand{\figOneScale}{1.0}

\usetikzlibrary{positioning,calc,fit,backgrounds,arrows,shapes.geometric}

% ---------------------------------------------------------------- palette
% Okabe & Ito (2008), CVD-safe qualitative set.
\providecommand{\oiOrangeDefined}{}
\definecolor{oiOrange}{HTML}{E69F00}
\definecolor{oiSky}   {HTML}{56B4E9}
\definecolor{oiGreen} {HTML}{009E73}
\definecolor{oiBlue}  {HTML}{0072B2}
\definecolor{oiPurple}{HTML}{CC79A7}
\definecolor{oiGrey}  {HTML}{7F7F7F}
\definecolor{oiGreyL} {HTML}{C4C4C4}
\definecolor{oiGreyD} {HTML}{4D4D4D}

% ---------------------------------------------------------------- type
% Body 10pt Times => \small 9pt, \footnotesize 8pt, \scriptsize 7pt.
% Titles are 9pt bold; all explanatory text is 8pt; 7pt is used only for
% short qualifiers.  Nothing in this figure is smaller than 7pt.
\tikzset{
  fpanel/.style   ={rounded corners=2.4pt, draw=#1!55!white, line width=0.7pt,
                    fill=#1!4},
  ftitle/.style   ={anchor=north west, font=\small\bfseries, inner sep=0pt,
                    text=oiGreyD},
  fsub/.style     ={anchor=north west, font=\footnotesize, inner sep=0pt,
                    text=oiGrey, align=left},
  fbody/.style    ={anchor=north west, font=\footnotesize, inner sep=0pt,
                    align=left},
  fnote/.style    ={anchor=north west, font=\scriptsize, inner sep=0pt,
                    align=left, text=oiGreyD},
  fctr/.style     ={anchor=north, font=\footnotesize, inner sep=0pt,
                    align=center},
  fctrn/.style    ={anchor=north, font=\scriptsize, inner sep=0pt,
                    align=center, text=oiGreyD},
  flow/.style     ={-stealth', line width=1.0pt, draw=oiBlue},
  teach/.style    ={-stealth', line width=1.0pt, draw=oiOrange!88!black},
  freearr/.style  ={stealth'-stealth', line width=0.8pt, draw=oiGrey,
                    dash pattern=on 2.4pt off 1.8pt},
  eqbox/.style    ={rounded corners=2.4pt, draw=oiGreen!65!white,
                    line width=0.8pt, fill=oiGreen!8},
}

% ------------------------------------------------------- molecule glyphs
% Deterministic coordinates: the figure is identical on every compile.
\newcommand{\fAtom}[3]{\fill[#3] (#1,#2) circle (2.05pt);%
  \draw[oiGreyD!55, line width=0.2pt] (#1,#2) circle (2.05pt);}
\newcommand{\fAtomS}[3]{\fill[#3] (#1,#2) circle (1.35pt);%
  \draw[oiGreyD!45, line width=0.15pt] (#1,#2) circle (1.35pt);}
\newcommand{\fBond}[4]{\draw[oiGrey!80, line width=0.7pt] (#1,#2) -- (#3,#4);}

% t = 0 : unstructured cloud, no bonds, no identities
\newcommand{\fMolNoise}{%
  \fAtomS{-0.36}{ 0.32}{oiGreyL}\fAtomS{ 0.12}{ 0.38}{oiGreyL}%
  \fAtomS{ 0.38}{ 0.06}{oiGreyL}\fAtomS{-0.10}{-0.02}{oiGreyL}%
  \fAtomS{-0.42}{-0.28}{oiGreyL}\fAtomS{ 0.06}{-0.36}{oiGreyL}%
  \fAtomS{ 0.42}{-0.34}{oiGreyL}\fAtomS{ 0.26}{ 0.24}{oiGreyL}%
  \fAtomS{-0.18}{ 0.12}{oiGreyL}}

% t = 1 : a structure with distinguishable element identities.  Element
% colours are neutral (grey / blue / sky / purple) so they cannot be read
% as the semantic green (value) or orange (teacher).
\newcommand{\fMolFinal}{%
  \fBond{ 0.000}{ 0.30}{-0.285}{ 0.09}\fBond{-0.285}{ 0.09}{-0.176}{-0.25}%
  \fBond{-0.176}{-0.25}{ 0.176}{-0.25}\fBond{ 0.176}{-0.25}{ 0.285}{ 0.09}%
  \fBond{ 0.285}{ 0.09}{ 0.000}{ 0.30}%
  \fBond{ 0.285}{ 0.09}{ 0.600}{ 0.27}%
  \fBond{-0.285}{ 0.09}{-0.560}{ 0.32}%
  \fBond{-0.176}{-0.25}{-0.385}{-0.47}%
  \fAtom{ 0.000}{ 0.30}{oiGreyD!75}\fAtom{-0.285}{ 0.09}{oiGreyD!75}%
  \fAtom{-0.176}{-0.25}{oiGreyD!75}\fAtom{ 0.176}{-0.25}{oiGreyD!75}%
  \fAtom{ 0.285}{ 0.09}{oiBlue}%
  \fAtom{ 0.600}{ 0.27}{oiPurple}%
  \fAtomS{-0.560}{ 0.32}{oiSky}%
  \fAtomS{-0.385}{-0.47}{white}}

% the same structure, small, for the displacement cloud in panel (b)
\newcommand{\fMolSmall}[1]{%
  \draw[#1, line width=0.55pt]
    ( 0.000, 0.24) -- (-0.228, 0.07) -- (-0.141,-0.20) --
    ( 0.141,-0.20) -- ( 0.228, 0.07) -- cycle;
  \draw[#1, line width=0.55pt] ( 0.228, 0.07) -- ( 0.470, 0.22);
  \draw[#1, line width=0.55pt] (-0.228, 0.07) -- (-0.440, 0.25);
  \foreach \px/\py in {0/0.24, -0.228/0.07, -0.141/-0.20, 0.141/-0.20,
                       0.228/0.07, 0.470/0.22, -0.440/0.25}
    {\fill[#1] (\px,\py) circle (1.15pt);}}

% =====================================================================
\begin{tikzpicture}[x=1cm, y=1cm, scale=\figOneScale, >=stealth']
% stable bounding box
\path (0,0) rectangle (13.90,8.80);

% ---- panel frames ----------------------------------------------------
\draw[fpanel=oiBlue]   (0.02,4.85) rectangle (6.82,8.75);
\draw[fpanel=oiOrange] (7.10,4.85) rectangle (13.88,8.75);
\draw[fpanel=oiGreen]  (0.02,0.05) rectangle (6.82,4.55);
\draw[fpanel=oiGrey]   (7.10,0.05) rectangle (13.88,4.55);

% =====================================================================
% (a) THE GENERATOR
% =====================================================================
\node[ftitle] at (0.22,8.66) {(a) A de novo flow};
\node[fsub, text width=6.30cm] at (0.22,8.30)
     {composition \emph{and} geometry from one model};

% prior
\begin{scope}[shift={(1.15,7.34)}]\fMolNoise\end{scope}
\node[fctr, text width=2.60cm] at (1.32,6.66)
     {$x_0\sim\mathcal{N}(0,\sigma^{2}I)$};
\node[fctrn, text width=2.40cm] at (1.32,6.28) {random types};

% flow arrow
\draw[flow] (2.20,7.34) -- (3.90,7.34);
\node[anchor=south, font=\footnotesize, text=oiBlue!80!black] at (3.05,7.46)
     {$v_\theta(x_t,t)$};
\node[anchor=north, font=\scriptsize, text=oiGreyD] at (3.05,7.20)
     {flow matching};

% sample
\begin{scope}[shift={(5.15,7.34)}]\fMolFinal\end{scope}
\node[fctr, text width=2.60cm] at (5.15,6.66) {$(c,\,x)$};
\node[fctrn, text width=3.10cm] at (5.15,6.30)
     {atom types $c$, coordinates $x$};

\node[fbody, text width=6.46cm] at (0.22,5.94)
     {No bond labels anywhere. With the discrete labels
      clamped, reversing the positional ODE gives
      $\log q_\theta(x\mid c)$.};

% =====================================================================
% (b) THE TEACHER
% =====================================================================
\node[ftitle] at (7.30,8.66) {(b) Teacher energy values};
\node[fsub, text width=6.30cm] at (7.30,8.30)
     {$K$ displaced geometries of \emph{one} parent};

% displacement cloud: two faint copies + one dark reference structure
\draw[oiGrey!45, line width=0.5pt, dash pattern=on 1.8pt off 1.6pt]
      (8.35,7.24) circle (0.72);
\begin{scope}[shift={(8.19,7.40)}]\fMolSmall{oiGreyL}\end{scope}
\begin{scope}[shift={(8.54,7.10)}]\fMolSmall{oiGreyL}\end{scope}
\begin{scope}[shift={(8.33,7.22)}]\fMolSmall{oiGreyD!80}\end{scope}
\node[fctr, text width=2.40cm] at (8.35,6.36) {$x_1,\dots,x_K$};
\node[fctrn, text width=2.60cm] at (8.35,6.04) {same composition $c$};

% teacher arrow
\draw[teach] (9.28,7.24) -- (10.22,7.24);
\node[anchor=south, font=\footnotesize, text=oiOrange!55!black] at (9.75,7.36)
     {$E(c,\cdot)$};

% energy ladder
\draw[-stealth', line width=0.7pt, draw=oiGreyD] (10.72,6.42) -- (10.72,7.94);
\node[anchor=east, font=\footnotesize, text=oiGreyD] at (10.62,7.84) {$E$};
\foreach \yy in {6.58, 6.86, 7.12, 7.44, 7.74}
  {\draw[oiGrey!70, line width=0.5pt] (10.72,\yy) -- (11.06,\yy);
   \fill[oiOrange!85!black] (11.06,\yy) circle (1.6pt);}
\node[anchor=west, font=\footnotesize] at (11.24,7.62)
     {$E_k=E(c,x_k)$};
\node[anchor=north west, font=\scriptsize, text width=2.50cm, text=oiGreyD]
      at (11.24,7.30) {one value per geometry};

\node[fbody, text width=6.34cm] at (7.30,5.62)
     {A universal neural potential over $83$ elements supplies every
      value. Only the values are used.};

% =====================================================================
% (c) THE OBJECTIVE
% =====================================================================
\node[ftitle] at (0.22,4.46) {(c) Proposed energy-value supervision};

\draw[eqbox] (0.16,2.86) rectangle (6.68,4.06);
\node[anchor=center, font=\small] at (3.42,3.86)
     {$\mathcal{L}_{\mathrm{ours}}
       =\mathcal{L}_{\mathrm{FM}}+\lambda_E\,\mathcal{L}_{\mathrm{value}}$};
\node[anchor=center, font=\small] at (3.42,3.48)
     {$\mathcal{L}_{\mathrm{value}}
       =\mathrm{Var}_{k}\big[\,\log q_\theta(x_k\mid c)+E_k/kT\,\big]$};
% The cancellation is the selling point of the objective, so it is drawn in
% the value colour rather than in neutral grey.  It is a statement about the
% objective, not a measured claim.
\node[anchor=north, font=\scriptsize, inner sep=0pt, align=center,
      text=oiGreen!45!black, text width=6.20cm] at (3.42,3.22)
     {one added term; the partition function cancels};

% schematic alignment plot
\draw[oiGreyD, line width=0.7pt] (1.36,2.52) -- (1.36,1.02) -- (3.98,1.02);
\node[rotate=90, anchor=south, font=\footnotesize, text=oiGreyD]
      at (1.20,1.77) {$\log q_\theta$};
\node[anchor=north, font=\footnotesize, text=oiGreyD] at (2.70,0.94)
     {$-E_k/kT$};
% the line the value term asks for: slope one, offset free
\draw[oiGreen!70!black, line width=0.9pt, dash pattern=on 2.6pt off 1.8pt]
      (1.52,1.16) -- (3.82,2.40);
% unconstrained (grey) and aligned (green) geometries
\foreach \px/\py in {1.68/2.12, 2.02/1.24, 2.36/2.32, 2.72/1.42,
                     3.08/1.98, 3.46/1.30, 3.72/2.20}
   {\fill[oiGrey!75] (\px,\py) circle (1.7pt);}
\foreach \px/\py in {1.68/1.22, 2.02/1.42, 2.36/1.50, 2.72/1.76,
                     3.08/1.92, 3.46/2.14, 3.72/2.36}
   {\fill[oiGreen] (\px,\py) circle (1.7pt);}
% legend
\fill[oiGrey!75] (4.42,2.32) circle (1.7pt);
\node[anchor=west, font=\footnotesize, text width=2.00cm] at (4.58,2.32)
     {$\mathcal{L}_{\mathrm{value}}$ large};
\fill[oiGreen] (4.42,1.86) circle (1.7pt);
\node[anchor=west, font=\footnotesize, text width=2.00cm] at (4.58,1.86)
     {$\mathcal{L}_{\mathrm{value}}=0$};
\node[fnote, text width=2.30cm] at (4.30,1.50)
     {one variance per group};

\node[anchor=north, font=\scriptsize, text=oiGreyD] at (3.42,0.56)
     {$\mathcal{L}_{\mathrm{value}}=0$ within a group
      $\iff \log q_\theta=-E_k/kT+b$};

% =====================================================================
% (d) THE SCOPE
% =====================================================================
% Two lines of \small\bfseries; the composition fact that used to sit in a
% separate subtitle is folded into the region label below, which is where the
% single conditional density is named anyway.
\node[ftitle, text width=6.50cm, align=left] at (7.30,4.46)
     {(d) Local supervision and the\\path to global control};

% the enclosing region: ONE conditional density, several supervised groups
\draw[rounded corners=2.8pt, draw=oiGrey!50, line width=0.6pt,
      dash pattern=on 2.2pt off 1.8pt] (7.42,1.02) rectangle (13.56,3.54);
\node[anchor=center, font=\scriptsize, text=oiGreyD, inner sep=1.6pt,
      fill=oiGrey!4] at (10.49,3.54)
     {one density $q_\theta(\cdot\mid c)$ --- same composition $c$};

\foreach \cx/\lbl in {8.86/A, 12.12/B}
{%
  \node[anchor=north, font=\footnotesize] at (\cx,3.36)
       {group $\lbl$ (basin $\lbl$)};
  \draw[oiGreen!60!white, line width=0.8pt, fill=oiGreen!7]
        (\cx,2.38) ellipse (1.28 and 0.56);
  \draw[oiGreen!45!black, line width=0.6pt]
        (\cx-0.88,2.26) -- (\cx-0.44,2.52) -- (\cx,2.36)
        -- (\cx+0.44,2.69) -- (\cx+0.84,2.50);
  \foreach \dx/\yy in {-0.88/2.26, -0.44/2.52, 0.00/2.36,
                        0.44/2.69, 0.84/2.50}
    {\fill[oiGreen] (\cx+\dx,\yy) circle (1.7pt);}
  \node[anchor=north, font=\scriptsize, text=oiGreyD] at (\cx,1.74)
       {local relation constrained};
  \node[anchor=north, font=\footnotesize] at (\cx,1.46)
       {offset $b_{\lbl}$ free};
}

\draw[freearr] (10.16,2.38) -- (10.82,2.38);
\node[anchor=south, font=\footnotesize, text=oiGrey] at (10.49,2.46) {?};

% Same mathematical content as before -- b_B - b_A and the relative mass are
% not identified -- with the constructive half of the statement added: what
% would have to be built to identify them.  No claim is weakened.
% WIDTH BUDGET.  The float sets \figOneScale = 1.000 and the picture is measured
% at that scale (see the measurement note in fig1_objective_float.tex: at 1.000 a
% 600 dpi render has zero ink outside the four panel frames; at 0.86 it has 4085
% such pixels, because tikz's `scale' shrinks COORDINATES but not type).  Panel
% (d) is 6.78cm wide at scale 1.000 and every font size here is absolute.  Two lines of at most ~42 \footnotesize characters is the whole
% budget; anything longer runs out of the panel and into panel (c).  Both qualifications survive the compression: the offset
% DIFFERENCE and the relative mass are the two unidentified quantities, and the
% second line names the extension that would identify them.
\node[anchor=north, font=\footnotesize, align=center]
      at (10.49,0.92)
     {$b_B-b_A$, relative mass: not identified.\\
      Connecting groups would share one offset.};

\end{tikzpicture}

\figOneFinish

  \caption{\textbf{Energy-value supervision without a partition function.}
    Ideal-objective schematic; no measured quantity appears here. The archived implementation does not realise this density construction (\Cref{sec:method:impl}).
    \textbf{(a)}~With the endpoint discrete labels clamped, reversing the
    positional ODE defines the \emph{clamped positional-flow density}
    $\log q_\theta(x \mid c)$, which need not equal the joint generator's
    conditional density. \textbf{(b)}~A universal neural potential values each
    of $K$ displaced geometries of one parent. \textbf{(c)}~The proposed
    objective is flow matching plus the within-group variance of
    $\log q_\theta(x_k \mid c) + E_k/kT$, which vanishes exactly when
    $\log q_\theta = -E_k/kT + b$ inside the group: the partition function
    cancels and the offset $b$ is free. \textbf{(d)}~Two disconnected groups of
    the \emph{same} density keep independent offsets, so the term constrains
    the local density--energy relation inside each group and leaves their
    relative mass unidentified; connecting groups so they share one offset is
    the extension toward global control. Panel detail in \Cref{app:figdetail}.}
  \label{fig:objective}
\end{figure}

%
%  The picture is drawn at 13.90cm = 5.47in, i.e. full ICLR single-column
%  width at scale 1.0.  Do NOT wrap it in \resizebox and do NOT scale it:
%  the in-figure type is \small / \footnotesize / \scriptsize of the 10pt
%  body font and scaling would desynchronise it from the caption.
%
%  PROVENANCE.  Conceptual figure, no data.  Built per the Figure 1 brief:
%  four panels (generator / teacher values / the value term / scope), and
%  the gradient-supervision branch of the superseded fig1_method.tex is
%  removed entirely rather than demoted.  No panel asserts a result.
%
%  REVISION C13 (figure-editor pass).  The framing moves to the method: panel
%  (c) is titled "Proposed energy-value supervision" and states inside the box
%  that the partition function cancels; panel (d)'s heading is "Local
%  supervision and the path to global control" and its closing line names the
%  extension (connecting groups so they share one offset) alongside the
%  unchanged qualification that b_B - b_A and the relative mass are not
%  identified.  No mathematical content was removed.
%
%  REVISION C8.  The caption is rewritten only where panels (c) and (d)
%  changed.  (c) now names the complete objective, so the caption states
%  it.  (d) now shows two disconnected groups of ONE composition, so the
%  caption no longer says the offset is free "between compositions"; the
%  freedom is already there inside a single conditional density.  The
%  verb "pins" is gone: the term CONSTRAINS THE LOCAL RELATION between
%  log-density and energy within a group.  The disconnection claim is
%  empirical and audited -- distinct parents sharing a composition supply
%  distinct perturbation clouds that share no sampled geometry (31% of
%  training parents sit on a repeated composition and are still
%  disconnected), so the overlap graph is a set of isolated vertices.
% =====================================================================
\begin{figure}[t]
  \centering
  \providecommand{\figOneScale}{1.0}\renewcommand{\figOneScale}{1.000}%
% SCALE = 1.000, THE FIGURE's NATURAL SIZE.  Do not lower it to buy page space.
% TikZ's `scale' multiplies COORDINATES but not type: every label in this
% picture is set at an absolute size (\small / \footnotesize / \scriptsize of
% the 10pt body), so any scale < 1 shrinks the panel frames out from under
% text that keeps its size, and the text runs out of the panels.
% MEASURED 2026-08-30 (600 dpi render of the standalone picture; panel frames
% located from the rendered strokes, so the calibration is self-checking at
% 236.22 px/cm on both axes):
%     scale 0.86 -> 4085 ink pixels outside the four panel frames, in 6
%                   regions.  Panel (b)'s body line runs to x = 14.59 cm,
%                   i.e. 0.71 cm past its own frame and 0.69 cm past the
%                   picture's 13.90 cm bounding box; panel (a)'s last two lines
%                   cross into panel (b)'s frame; panel (c)'s closing identity
%                   line crosses the left frame at x = -0.17 cm.
%     scale 1.000 -> 0 ink pixels outside the panels, at every ink threshold
%                   tested (lum < 128 and lum < 200, with 5 px of stroke
%                   tolerance).  All four panels are clean.
% PAGE COST: none that changes the budget.  With 1.000 the main text ends on
% page 9 and a probe \vspace inserted before \label{endofmaintext} shows a
% further +108pt fits before it spills to page 10 (+114pt spills); at 0.86 the
% same probe breaks between +108pt and +120pt.  The fix costs at most about
% one body line of final-page headroom.
  % =====================================================================
%  fig1_objective.tex --- Figure 1: the objective and its scope
%
%  CONCEPTUAL FIGURE.  NO DATA.  No number in this figure is measured;
%  every glyph is a schematic.  Nothing here asserts a result.
%
%  Replaces the old five-band figures/fig1_method.tex, which was too
%  dense for a first-page explanatory figure and whose bands (b)/(c)
%  gave the gradient-supervision path equal visual status.  Per the
%  build brief the force branch is REMOVED from this figure entirely;
%  gradient supervision appears only later, as an experimental
%  comparator, and is not drawn here.
%
%  Pure TikZ.  No external images, no pgfplots, no \includegraphics,
%  no arrows.meta, no pgfmathrandom (byte-reproducible across compiles).
%  Okabe--Ito colour-blind-safe palette, semantics shared with
%  make_experiment_figures.py::PALETTE:
%     blue    model / flow / learned quantities
%     green   the value term (energy values, L_value)
%     orange  the teacher (universal neural potential)
%     grey    neutral structure
%  Vermilion (= gradient supervision) and purple (= scrambled control)
%  are deliberately ABSENT: neither concept appears in this figure.
%
%  USAGE
%   (a) standalone (verification / slides):   pdflatex fig1_objective.tex
%   (b) inside the paper -- main.tex already \usepackage{tikz} and
%       \usepackage{amsmath}:
%         % =====================================================================
%  fig1_objective_float.tex --- the float wrapper for Figure 1.
%
%  Figure 1 is pure TikZ, so this wrapper \input's the picture directly
%  rather than \includegraphics-ing a PDF (the same mechanism the old
%  fig1_method.tex used).  main.tex already loads tikz and amsmath.
%
%  Include from sections/01_intro.tex with:
%      \input{figures/fig1_objective_float}
%
%  The picture is drawn at 13.90cm = 5.47in, i.e. full ICLR single-column
%  width at scale 1.0.  Do NOT wrap it in \resizebox and do NOT scale it:
%  the in-figure type is \small / \footnotesize / \scriptsize of the 10pt
%  body font and scaling would desynchronise it from the caption.
%
%  PROVENANCE.  Conceptual figure, no data.  Built per the Figure 1 brief:
%  four panels (generator / teacher values / the value term / scope), and
%  the gradient-supervision branch of the superseded fig1_method.tex is
%  removed entirely rather than demoted.  No panel asserts a result.
%
%  REVISION C13 (figure-editor pass).  The framing moves to the method: panel
%  (c) is titled "Proposed energy-value supervision" and states inside the box
%  that the partition function cancels; panel (d)'s heading is "Local
%  supervision and the path to global control" and its closing line names the
%  extension (connecting groups so they share one offset) alongside the
%  unchanged qualification that b_B - b_A and the relative mass are not
%  identified.  No mathematical content was removed.
%
%  REVISION C8.  The caption is rewritten only where panels (c) and (d)
%  changed.  (c) now names the complete objective, so the caption states
%  it.  (d) now shows two disconnected groups of ONE composition, so the
%  caption no longer says the offset is free "between compositions"; the
%  freedom is already there inside a single conditional density.  The
%  verb "pins" is gone: the term CONSTRAINS THE LOCAL RELATION between
%  log-density and energy within a group.  The disconnection claim is
%  empirical and audited -- distinct parents sharing a composition supply
%  distinct perturbation clouds that share no sampled geometry (31% of
%  training parents sit on a repeated composition and are still
%  disconnected), so the overlap graph is a set of isolated vertices.
% =====================================================================
\begin{figure}[t]
  \centering
  \providecommand{\figOneScale}{1.0}\renewcommand{\figOneScale}{1.000}%
% SCALE = 1.000, THE FIGURE's NATURAL SIZE.  Do not lower it to buy page space.
% TikZ's `scale' multiplies COORDINATES but not type: every label in this
% picture is set at an absolute size (\small / \footnotesize / \scriptsize of
% the 10pt body), so any scale < 1 shrinks the panel frames out from under
% text that keeps its size, and the text runs out of the panels.
% MEASURED 2026-08-30 (600 dpi render of the standalone picture; panel frames
% located from the rendered strokes, so the calibration is self-checking at
% 236.22 px/cm on both axes):
%     scale 0.86 -> 4085 ink pixels outside the four panel frames, in 6
%                   regions.  Panel (b)'s body line runs to x = 14.59 cm,
%                   i.e. 0.71 cm past its own frame and 0.69 cm past the
%                   picture's 13.90 cm bounding box; panel (a)'s last two lines
%                   cross into panel (b)'s frame; panel (c)'s closing identity
%                   line crosses the left frame at x = -0.17 cm.
%     scale 1.000 -> 0 ink pixels outside the panels, at every ink threshold
%                   tested (lum < 128 and lum < 200, with 5 px of stroke
%                   tolerance).  All four panels are clean.
% PAGE COST: none that changes the budget.  With 1.000 the main text ends on
% page 9 and a probe \vspace inserted before \label{endofmaintext} shows a
% further +108pt fits before it spills to page 10 (+114pt spills); at 0.86 the
% same probe breaks between +108pt and +120pt.  The fix costs at most about
% one body line of final-page headroom.
  \input{figures/fig1_objective}
  \caption{\textbf{Energy-value supervision without a partition function.}
    Ideal-objective schematic; no measured quantity appears here. The archived implementation does not realise this density construction (\Cref{sec:method:impl}).
    \textbf{(a)}~With the endpoint discrete labels clamped, reversing the
    positional ODE defines the \emph{clamped positional-flow density}
    $\log q_\theta(x \mid c)$, which need not equal the joint generator's
    conditional density. \textbf{(b)}~A universal neural potential values each
    of $K$ displaced geometries of one parent. \textbf{(c)}~The proposed
    objective is flow matching plus the within-group variance of
    $\log q_\theta(x_k \mid c) + E_k/kT$, which vanishes exactly when
    $\log q_\theta = -E_k/kT + b$ inside the group: the partition function
    cancels and the offset $b$ is free. \textbf{(d)}~Two disconnected groups of
    the \emph{same} density keep independent offsets, so the term constrains
    the local density--energy relation inside each group and leaves their
    relative mass unidentified; connecting groups so they share one offset is
    the extension toward global control. Panel detail in \Cref{app:figdetail}.}
  \label{fig:objective}
\end{figure}

%       or, if the writer prefers to own the float:
%         \begin{figure}[t]\centering
%           % =====================================================================
%  fig1_objective.tex --- Figure 1: the objective and its scope
%
%  CONCEPTUAL FIGURE.  NO DATA.  No number in this figure is measured;
%  every glyph is a schematic.  Nothing here asserts a result.
%
%  Replaces the old five-band figures/fig1_method.tex, which was too
%  dense for a first-page explanatory figure and whose bands (b)/(c)
%  gave the gradient-supervision path equal visual status.  Per the
%  build brief the force branch is REMOVED from this figure entirely;
%  gradient supervision appears only later, as an experimental
%  comparator, and is not drawn here.
%
%  Pure TikZ.  No external images, no pgfplots, no \includegraphics,
%  no arrows.meta, no pgfmathrandom (byte-reproducible across compiles).
%  Okabe--Ito colour-blind-safe palette, semantics shared with
%  make_experiment_figures.py::PALETTE:
%     blue    model / flow / learned quantities
%     green   the value term (energy values, L_value)
%     orange  the teacher (universal neural potential)
%     grey    neutral structure
%  Vermilion (= gradient supervision) and purple (= scrambled control)
%  are deliberately ABSENT: neither concept appears in this figure.
%
%  USAGE
%   (a) standalone (verification / slides):   pdflatex fig1_objective.tex
%   (b) inside the paper -- main.tex already \usepackage{tikz} and
%       \usepackage{amsmath}:
%         \input{figures/fig1_objective_float}
%       or, if the writer prefers to own the float:
%         \begin{figure}[t]\centering
%           \input{figures/fig1_objective}
%           \caption{...}\label{fig:objective}
%         \end{figure}
%
%  SIZE.  Natural size 13.90cm x 8.80cm == 5.47in x 3.46in.
%  ICLR is a SINGLE 5.5in text column, so this is a full-text-width
%  figure at scale 1.0.  DO NOT rescale it and DO NOT wrap it in
%  \resizebox: every glyph inside is \small / \footnotesize /
%  \scriptsize of the 10pt body font, and scaling desynchronises
%  in-figure text from the caption.
%
%  LAYOUT.  A 2x2 grid, read left-to-right, top-to-bottom:
%    (a) the generator      composition AND geometry from one flow
%    (b) the teacher        energy values for K displacements of one parent
%    (c) the objective      L_ours = L_FM + lambda_E L_value, and nothing else
%    (d) the scope          two basins of ONE composition; offset free
%  Panel widths 6.80 / 6.78 cm; no wire crosses a panel boundary, so
%  there is no routing to collide with text.
%
%  REVISION C8 (this pass).  Two changes, both confined to panels (c)/(d):
%   * (c) now displays the complete objective, L_ours = L_FM + lambda_E
%     L_value, so the reader cannot infer a hidden force or anchor branch.
%     Neither branch was ever drawn here and neither is added.
%   * (d) previously drew two groups of two DIFFERENT compositions
%     (c_A, c_B).  That was wrong twice over.  It misstated the failure
%     mode -- the unidentified offset is already present WITHIN a single
%     conditional density -- and it flirted with the claim that one
%     composition contributes exactly one group, which the composition
%     audit refutes: in the training perturbation shard
%     (processed_data/omol25_4m_processed/perturbation_train_n30000_s0.pt)
%     30,000 parent groups carry only
%     22,822 distinct (element multiset, total charge) keys, 2,204 keys
%     are shared by more than one parent (9,382 parents, 31%; max
%     multiplicity 80).  Panel (d) therefore now shows two groups A and B
%     of the SAME composition c: two disconnected supervised regions of
%     one density p_theta(.|c), each with its own free offset b_A, b_B.
%     "ordering fixed" / "pinned" wording is replaced throughout by
%     "the local relation between log-density and energy is constrained".
%  In the EVALUATION population (val_data_processed.pt; the 120 grid
%  parents and the 93 primary parents) every parent has a distinct
%  composition, multiplicity 1 -- that is an audited property of the
%  evaluation set, not of the supervision design, and this figure does
%  not assert it.
% =====================================================================

% ---- standalone bootstrap (skipped when \input into a tikz-enabled doc) ----
\ifx\tikzpicture\undefined
  \documentclass[10pt]{article}
  % 0.05cm of slack past the 13.90cm bounding box so that the panel border
  % stroke width does not report a spurious overfull \hbox in standalone mode.
  \usepackage[paperwidth=14.40cm,paperheight=9.20cm,margin=0.20cm,noheadfoot]{geometry}
  \usepackage{amsmath,amssymb,bm}
  \usepackage{times}
  \usepackage{tikz}
  \setlength{\topskip}{0pt}
  \setlength{\parindent}{0pt}
  \pagestyle{empty}
  \begin{document}%
  \def\figOneFinish{\end{document}}
\fi
\providecommand{\figOneFinish}{\relax}
\providecommand{\figOneScale}{1.0}

\usetikzlibrary{positioning,calc,fit,backgrounds,arrows,shapes.geometric}

% ---------------------------------------------------------------- palette
% Okabe & Ito (2008), CVD-safe qualitative set.
\providecommand{\oiOrangeDefined}{}
\definecolor{oiOrange}{HTML}{E69F00}
\definecolor{oiSky}   {HTML}{56B4E9}
\definecolor{oiGreen} {HTML}{009E73}
\definecolor{oiBlue}  {HTML}{0072B2}
\definecolor{oiPurple}{HTML}{CC79A7}
\definecolor{oiGrey}  {HTML}{7F7F7F}
\definecolor{oiGreyL} {HTML}{C4C4C4}
\definecolor{oiGreyD} {HTML}{4D4D4D}

% ---------------------------------------------------------------- type
% Body 10pt Times => \small 9pt, \footnotesize 8pt, \scriptsize 7pt.
% Titles are 9pt bold; all explanatory text is 8pt; 7pt is used only for
% short qualifiers.  Nothing in this figure is smaller than 7pt.
\tikzset{
  fpanel/.style   ={rounded corners=2.4pt, draw=#1!55!white, line width=0.7pt,
                    fill=#1!4},
  ftitle/.style   ={anchor=north west, font=\small\bfseries, inner sep=0pt,
                    text=oiGreyD},
  fsub/.style     ={anchor=north west, font=\footnotesize, inner sep=0pt,
                    text=oiGrey, align=left},
  fbody/.style    ={anchor=north west, font=\footnotesize, inner sep=0pt,
                    align=left},
  fnote/.style    ={anchor=north west, font=\scriptsize, inner sep=0pt,
                    align=left, text=oiGreyD},
  fctr/.style     ={anchor=north, font=\footnotesize, inner sep=0pt,
                    align=center},
  fctrn/.style    ={anchor=north, font=\scriptsize, inner sep=0pt,
                    align=center, text=oiGreyD},
  flow/.style     ={-stealth', line width=1.0pt, draw=oiBlue},
  teach/.style    ={-stealth', line width=1.0pt, draw=oiOrange!88!black},
  freearr/.style  ={stealth'-stealth', line width=0.8pt, draw=oiGrey,
                    dash pattern=on 2.4pt off 1.8pt},
  eqbox/.style    ={rounded corners=2.4pt, draw=oiGreen!65!white,
                    line width=0.8pt, fill=oiGreen!8},
}

% ------------------------------------------------------- molecule glyphs
% Deterministic coordinates: the figure is identical on every compile.
\newcommand{\fAtom}[3]{\fill[#3] (#1,#2) circle (2.05pt);%
  \draw[oiGreyD!55, line width=0.2pt] (#1,#2) circle (2.05pt);}
\newcommand{\fAtomS}[3]{\fill[#3] (#1,#2) circle (1.35pt);%
  \draw[oiGreyD!45, line width=0.15pt] (#1,#2) circle (1.35pt);}
\newcommand{\fBond}[4]{\draw[oiGrey!80, line width=0.7pt] (#1,#2) -- (#3,#4);}

% t = 0 : unstructured cloud, no bonds, no identities
\newcommand{\fMolNoise}{%
  \fAtomS{-0.36}{ 0.32}{oiGreyL}\fAtomS{ 0.12}{ 0.38}{oiGreyL}%
  \fAtomS{ 0.38}{ 0.06}{oiGreyL}\fAtomS{-0.10}{-0.02}{oiGreyL}%
  \fAtomS{-0.42}{-0.28}{oiGreyL}\fAtomS{ 0.06}{-0.36}{oiGreyL}%
  \fAtomS{ 0.42}{-0.34}{oiGreyL}\fAtomS{ 0.26}{ 0.24}{oiGreyL}%
  \fAtomS{-0.18}{ 0.12}{oiGreyL}}

% t = 1 : a structure with distinguishable element identities.  Element
% colours are neutral (grey / blue / sky / purple) so they cannot be read
% as the semantic green (value) or orange (teacher).
\newcommand{\fMolFinal}{%
  \fBond{ 0.000}{ 0.30}{-0.285}{ 0.09}\fBond{-0.285}{ 0.09}{-0.176}{-0.25}%
  \fBond{-0.176}{-0.25}{ 0.176}{-0.25}\fBond{ 0.176}{-0.25}{ 0.285}{ 0.09}%
  \fBond{ 0.285}{ 0.09}{ 0.000}{ 0.30}%
  \fBond{ 0.285}{ 0.09}{ 0.600}{ 0.27}%
  \fBond{-0.285}{ 0.09}{-0.560}{ 0.32}%
  \fBond{-0.176}{-0.25}{-0.385}{-0.47}%
  \fAtom{ 0.000}{ 0.30}{oiGreyD!75}\fAtom{-0.285}{ 0.09}{oiGreyD!75}%
  \fAtom{-0.176}{-0.25}{oiGreyD!75}\fAtom{ 0.176}{-0.25}{oiGreyD!75}%
  \fAtom{ 0.285}{ 0.09}{oiBlue}%
  \fAtom{ 0.600}{ 0.27}{oiPurple}%
  \fAtomS{-0.560}{ 0.32}{oiSky}%
  \fAtomS{-0.385}{-0.47}{white}}

% the same structure, small, for the displacement cloud in panel (b)
\newcommand{\fMolSmall}[1]{%
  \draw[#1, line width=0.55pt]
    ( 0.000, 0.24) -- (-0.228, 0.07) -- (-0.141,-0.20) --
    ( 0.141,-0.20) -- ( 0.228, 0.07) -- cycle;
  \draw[#1, line width=0.55pt] ( 0.228, 0.07) -- ( 0.470, 0.22);
  \draw[#1, line width=0.55pt] (-0.228, 0.07) -- (-0.440, 0.25);
  \foreach \px/\py in {0/0.24, -0.228/0.07, -0.141/-0.20, 0.141/-0.20,
                       0.228/0.07, 0.470/0.22, -0.440/0.25}
    {\fill[#1] (\px,\py) circle (1.15pt);}}

% =====================================================================
\begin{tikzpicture}[x=1cm, y=1cm, scale=\figOneScale, >=stealth']
% stable bounding box
\path (0,0) rectangle (13.90,8.80);

% ---- panel frames ----------------------------------------------------
\draw[fpanel=oiBlue]   (0.02,4.85) rectangle (6.82,8.75);
\draw[fpanel=oiOrange] (7.10,4.85) rectangle (13.88,8.75);
\draw[fpanel=oiGreen]  (0.02,0.05) rectangle (6.82,4.55);
\draw[fpanel=oiGrey]   (7.10,0.05) rectangle (13.88,4.55);

% =====================================================================
% (a) THE GENERATOR
% =====================================================================
\node[ftitle] at (0.22,8.66) {(a) A de novo flow};
\node[fsub, text width=6.30cm] at (0.22,8.30)
     {composition \emph{and} geometry from one model};

% prior
\begin{scope}[shift={(1.15,7.34)}]\fMolNoise\end{scope}
\node[fctr, text width=2.60cm] at (1.32,6.66)
     {$x_0\sim\mathcal{N}(0,\sigma^{2}I)$};
\node[fctrn, text width=2.40cm] at (1.32,6.28) {random types};

% flow arrow
\draw[flow] (2.20,7.34) -- (3.90,7.34);
\node[anchor=south, font=\footnotesize, text=oiBlue!80!black] at (3.05,7.46)
     {$v_\theta(x_t,t)$};
\node[anchor=north, font=\scriptsize, text=oiGreyD] at (3.05,7.20)
     {flow matching};

% sample
\begin{scope}[shift={(5.15,7.34)}]\fMolFinal\end{scope}
\node[fctr, text width=2.60cm] at (5.15,6.66) {$(c,\,x)$};
\node[fctrn, text width=3.10cm] at (5.15,6.30)
     {atom types $c$, coordinates $x$};

\node[fbody, text width=6.46cm] at (0.22,5.94)
     {No bond labels anywhere. With the discrete labels
      clamped, reversing the positional ODE gives
      $\log q_\theta(x\mid c)$.};

% =====================================================================
% (b) THE TEACHER
% =====================================================================
\node[ftitle] at (7.30,8.66) {(b) Teacher energy values};
\node[fsub, text width=6.30cm] at (7.30,8.30)
     {$K$ displaced geometries of \emph{one} parent};

% displacement cloud: two faint copies + one dark reference structure
\draw[oiGrey!45, line width=0.5pt, dash pattern=on 1.8pt off 1.6pt]
      (8.35,7.24) circle (0.72);
\begin{scope}[shift={(8.19,7.40)}]\fMolSmall{oiGreyL}\end{scope}
\begin{scope}[shift={(8.54,7.10)}]\fMolSmall{oiGreyL}\end{scope}
\begin{scope}[shift={(8.33,7.22)}]\fMolSmall{oiGreyD!80}\end{scope}
\node[fctr, text width=2.40cm] at (8.35,6.36) {$x_1,\dots,x_K$};
\node[fctrn, text width=2.60cm] at (8.35,6.04) {same composition $c$};

% teacher arrow
\draw[teach] (9.28,7.24) -- (10.22,7.24);
\node[anchor=south, font=\footnotesize, text=oiOrange!55!black] at (9.75,7.36)
     {$E(c,\cdot)$};

% energy ladder
\draw[-stealth', line width=0.7pt, draw=oiGreyD] (10.72,6.42) -- (10.72,7.94);
\node[anchor=east, font=\footnotesize, text=oiGreyD] at (10.62,7.84) {$E$};
\foreach \yy in {6.58, 6.86, 7.12, 7.44, 7.74}
  {\draw[oiGrey!70, line width=0.5pt] (10.72,\yy) -- (11.06,\yy);
   \fill[oiOrange!85!black] (11.06,\yy) circle (1.6pt);}
\node[anchor=west, font=\footnotesize] at (11.24,7.62)
     {$E_k=E(c,x_k)$};
\node[anchor=north west, font=\scriptsize, text width=2.50cm, text=oiGreyD]
      at (11.24,7.30) {one value per geometry};

\node[fbody, text width=6.34cm] at (7.30,5.62)
     {A universal neural potential over $83$ elements supplies every
      value. Only the values are used.};

% =====================================================================
% (c) THE OBJECTIVE
% =====================================================================
\node[ftitle] at (0.22,4.46) {(c) Proposed energy-value supervision};

\draw[eqbox] (0.16,2.86) rectangle (6.68,4.06);
\node[anchor=center, font=\small] at (3.42,3.86)
     {$\mathcal{L}_{\mathrm{ours}}
       =\mathcal{L}_{\mathrm{FM}}+\lambda_E\,\mathcal{L}_{\mathrm{value}}$};
\node[anchor=center, font=\small] at (3.42,3.48)
     {$\mathcal{L}_{\mathrm{value}}
       =\mathrm{Var}_{k}\big[\,\log q_\theta(x_k\mid c)+E_k/kT\,\big]$};
% The cancellation is the selling point of the objective, so it is drawn in
% the value colour rather than in neutral grey.  It is a statement about the
% objective, not a measured claim.
\node[anchor=north, font=\scriptsize, inner sep=0pt, align=center,
      text=oiGreen!45!black, text width=6.20cm] at (3.42,3.22)
     {one added term; the partition function cancels};

% schematic alignment plot
\draw[oiGreyD, line width=0.7pt] (1.36,2.52) -- (1.36,1.02) -- (3.98,1.02);
\node[rotate=90, anchor=south, font=\footnotesize, text=oiGreyD]
      at (1.20,1.77) {$\log q_\theta$};
\node[anchor=north, font=\footnotesize, text=oiGreyD] at (2.70,0.94)
     {$-E_k/kT$};
% the line the value term asks for: slope one, offset free
\draw[oiGreen!70!black, line width=0.9pt, dash pattern=on 2.6pt off 1.8pt]
      (1.52,1.16) -- (3.82,2.40);
% unconstrained (grey) and aligned (green) geometries
\foreach \px/\py in {1.68/2.12, 2.02/1.24, 2.36/2.32, 2.72/1.42,
                     3.08/1.98, 3.46/1.30, 3.72/2.20}
   {\fill[oiGrey!75] (\px,\py) circle (1.7pt);}
\foreach \px/\py in {1.68/1.22, 2.02/1.42, 2.36/1.50, 2.72/1.76,
                     3.08/1.92, 3.46/2.14, 3.72/2.36}
   {\fill[oiGreen] (\px,\py) circle (1.7pt);}
% legend
\fill[oiGrey!75] (4.42,2.32) circle (1.7pt);
\node[anchor=west, font=\footnotesize, text width=2.00cm] at (4.58,2.32)
     {$\mathcal{L}_{\mathrm{value}}$ large};
\fill[oiGreen] (4.42,1.86) circle (1.7pt);
\node[anchor=west, font=\footnotesize, text width=2.00cm] at (4.58,1.86)
     {$\mathcal{L}_{\mathrm{value}}=0$};
\node[fnote, text width=2.30cm] at (4.30,1.50)
     {one variance per group};

\node[anchor=north, font=\scriptsize, text=oiGreyD] at (3.42,0.56)
     {$\mathcal{L}_{\mathrm{value}}=0$ within a group
      $\iff \log q_\theta=-E_k/kT+b$};

% =====================================================================
% (d) THE SCOPE
% =====================================================================
% Two lines of \small\bfseries; the composition fact that used to sit in a
% separate subtitle is folded into the region label below, which is where the
% single conditional density is named anyway.
\node[ftitle, text width=6.50cm, align=left] at (7.30,4.46)
     {(d) Local supervision and the\\path to global control};

% the enclosing region: ONE conditional density, several supervised groups
\draw[rounded corners=2.8pt, draw=oiGrey!50, line width=0.6pt,
      dash pattern=on 2.2pt off 1.8pt] (7.42,1.02) rectangle (13.56,3.54);
\node[anchor=center, font=\scriptsize, text=oiGreyD, inner sep=1.6pt,
      fill=oiGrey!4] at (10.49,3.54)
     {one density $q_\theta(\cdot\mid c)$ --- same composition $c$};

\foreach \cx/\lbl in {8.86/A, 12.12/B}
{%
  \node[anchor=north, font=\footnotesize] at (\cx,3.36)
       {group $\lbl$ (basin $\lbl$)};
  \draw[oiGreen!60!white, line width=0.8pt, fill=oiGreen!7]
        (\cx,2.38) ellipse (1.28 and 0.56);
  \draw[oiGreen!45!black, line width=0.6pt]
        (\cx-0.88,2.26) -- (\cx-0.44,2.52) -- (\cx,2.36)
        -- (\cx+0.44,2.69) -- (\cx+0.84,2.50);
  \foreach \dx/\yy in {-0.88/2.26, -0.44/2.52, 0.00/2.36,
                        0.44/2.69, 0.84/2.50}
    {\fill[oiGreen] (\cx+\dx,\yy) circle (1.7pt);}
  \node[anchor=north, font=\scriptsize, text=oiGreyD] at (\cx,1.74)
       {local relation constrained};
  \node[anchor=north, font=\footnotesize] at (\cx,1.46)
       {offset $b_{\lbl}$ free};
}

\draw[freearr] (10.16,2.38) -- (10.82,2.38);
\node[anchor=south, font=\footnotesize, text=oiGrey] at (10.49,2.46) {?};

% Same mathematical content as before -- b_B - b_A and the relative mass are
% not identified -- with the constructive half of the statement added: what
% would have to be built to identify them.  No claim is weakened.
% WIDTH BUDGET.  The float sets \figOneScale = 1.000 and the picture is measured
% at that scale (see the measurement note in fig1_objective_float.tex: at 1.000 a
% 600 dpi render has zero ink outside the four panel frames; at 0.86 it has 4085
% such pixels, because tikz's `scale' shrinks COORDINATES but not type).  Panel
% (d) is 6.78cm wide at scale 1.000 and every font size here is absolute.  Two lines of at most ~42 \footnotesize characters is the whole
% budget; anything longer runs out of the panel and into panel (c).  Both qualifications survive the compression: the offset
% DIFFERENCE and the relative mass are the two unidentified quantities, and the
% second line names the extension that would identify them.
\node[anchor=north, font=\footnotesize, align=center]
      at (10.49,0.92)
     {$b_B-b_A$, relative mass: not identified.\\
      Connecting groups would share one offset.};

\end{tikzpicture}

\figOneFinish

%           \caption{...}\label{fig:objective}
%         \end{figure}
%
%  SIZE.  Natural size 13.90cm x 8.80cm == 5.47in x 3.46in.
%  ICLR is a SINGLE 5.5in text column, so this is a full-text-width
%  figure at scale 1.0.  DO NOT rescale it and DO NOT wrap it in
%  \resizebox: every glyph inside is \small / \footnotesize /
%  \scriptsize of the 10pt body font, and scaling desynchronises
%  in-figure text from the caption.
%
%  LAYOUT.  A 2x2 grid, read left-to-right, top-to-bottom:
%    (a) the generator      composition AND geometry from one flow
%    (b) the teacher        energy values for K displacements of one parent
%    (c) the objective      L_ours = L_FM + lambda_E L_value, and nothing else
%    (d) the scope          two basins of ONE composition; offset free
%  Panel widths 6.80 / 6.78 cm; no wire crosses a panel boundary, so
%  there is no routing to collide with text.
%
%  REVISION C8 (this pass).  Two changes, both confined to panels (c)/(d):
%   * (c) now displays the complete objective, L_ours = L_FM + lambda_E
%     L_value, so the reader cannot infer a hidden force or anchor branch.
%     Neither branch was ever drawn here and neither is added.
%   * (d) previously drew two groups of two DIFFERENT compositions
%     (c_A, c_B).  That was wrong twice over.  It misstated the failure
%     mode -- the unidentified offset is already present WITHIN a single
%     conditional density -- and it flirted with the claim that one
%     composition contributes exactly one group, which the composition
%     audit refutes: in the training perturbation shard
%     (processed_data/omol25_4m_processed/perturbation_train_n30000_s0.pt)
%     30,000 parent groups carry only
%     22,822 distinct (element multiset, total charge) keys, 2,204 keys
%     are shared by more than one parent (9,382 parents, 31%; max
%     multiplicity 80).  Panel (d) therefore now shows two groups A and B
%     of the SAME composition c: two disconnected supervised regions of
%     one density p_theta(.|c), each with its own free offset b_A, b_B.
%     "ordering fixed" / "pinned" wording is replaced throughout by
%     "the local relation between log-density and energy is constrained".
%  In the EVALUATION population (val_data_processed.pt; the 120 grid
%  parents and the 93 primary parents) every parent has a distinct
%  composition, multiplicity 1 -- that is an audited property of the
%  evaluation set, not of the supervision design, and this figure does
%  not assert it.
% =====================================================================

% ---- standalone bootstrap (skipped when \input into a tikz-enabled doc) ----
\ifx\tikzpicture\undefined
  \documentclass[10pt]{article}
  % 0.05cm of slack past the 13.90cm bounding box so that the panel border
  % stroke width does not report a spurious overfull \hbox in standalone mode.
  \usepackage[paperwidth=14.40cm,paperheight=9.20cm,margin=0.20cm,noheadfoot]{geometry}
  \usepackage{amsmath,amssymb,bm}
  \usepackage{times}
  \usepackage{tikz}
  \setlength{\topskip}{0pt}
  \setlength{\parindent}{0pt}
  \pagestyle{empty}
  \begin{document}%
  \def\figOneFinish{\end{document}}
\fi
\providecommand{\figOneFinish}{\relax}
\providecommand{\figOneScale}{1.0}

\usetikzlibrary{positioning,calc,fit,backgrounds,arrows,shapes.geometric}

% ---------------------------------------------------------------- palette
% Okabe & Ito (2008), CVD-safe qualitative set.
\providecommand{\oiOrangeDefined}{}
\definecolor{oiOrange}{HTML}{E69F00}
\definecolor{oiSky}   {HTML}{56B4E9}
\definecolor{oiGreen} {HTML}{009E73}
\definecolor{oiBlue}  {HTML}{0072B2}
\definecolor{oiPurple}{HTML}{CC79A7}
\definecolor{oiGrey}  {HTML}{7F7F7F}
\definecolor{oiGreyL} {HTML}{C4C4C4}
\definecolor{oiGreyD} {HTML}{4D4D4D}

% ---------------------------------------------------------------- type
% Body 10pt Times => \small 9pt, \footnotesize 8pt, \scriptsize 7pt.
% Titles are 9pt bold; all explanatory text is 8pt; 7pt is used only for
% short qualifiers.  Nothing in this figure is smaller than 7pt.
\tikzset{
  fpanel/.style   ={rounded corners=2.4pt, draw=#1!55!white, line width=0.7pt,
                    fill=#1!4},
  ftitle/.style   ={anchor=north west, font=\small\bfseries, inner sep=0pt,
                    text=oiGreyD},
  fsub/.style     ={anchor=north west, font=\footnotesize, inner sep=0pt,
                    text=oiGrey, align=left},
  fbody/.style    ={anchor=north west, font=\footnotesize, inner sep=0pt,
                    align=left},
  fnote/.style    ={anchor=north west, font=\scriptsize, inner sep=0pt,
                    align=left, text=oiGreyD},
  fctr/.style     ={anchor=north, font=\footnotesize, inner sep=0pt,
                    align=center},
  fctrn/.style    ={anchor=north, font=\scriptsize, inner sep=0pt,
                    align=center, text=oiGreyD},
  flow/.style     ={-stealth', line width=1.0pt, draw=oiBlue},
  teach/.style    ={-stealth', line width=1.0pt, draw=oiOrange!88!black},
  freearr/.style  ={stealth'-stealth', line width=0.8pt, draw=oiGrey,
                    dash pattern=on 2.4pt off 1.8pt},
  eqbox/.style    ={rounded corners=2.4pt, draw=oiGreen!65!white,
                    line width=0.8pt, fill=oiGreen!8},
}

% ------------------------------------------------------- molecule glyphs
% Deterministic coordinates: the figure is identical on every compile.
\newcommand{\fAtom}[3]{\fill[#3] (#1,#2) circle (2.05pt);%
  \draw[oiGreyD!55, line width=0.2pt] (#1,#2) circle (2.05pt);}
\newcommand{\fAtomS}[3]{\fill[#3] (#1,#2) circle (1.35pt);%
  \draw[oiGreyD!45, line width=0.15pt] (#1,#2) circle (1.35pt);}
\newcommand{\fBond}[4]{\draw[oiGrey!80, line width=0.7pt] (#1,#2) -- (#3,#4);}

% t = 0 : unstructured cloud, no bonds, no identities
\newcommand{\fMolNoise}{%
  \fAtomS{-0.36}{ 0.32}{oiGreyL}\fAtomS{ 0.12}{ 0.38}{oiGreyL}%
  \fAtomS{ 0.38}{ 0.06}{oiGreyL}\fAtomS{-0.10}{-0.02}{oiGreyL}%
  \fAtomS{-0.42}{-0.28}{oiGreyL}\fAtomS{ 0.06}{-0.36}{oiGreyL}%
  \fAtomS{ 0.42}{-0.34}{oiGreyL}\fAtomS{ 0.26}{ 0.24}{oiGreyL}%
  \fAtomS{-0.18}{ 0.12}{oiGreyL}}

% t = 1 : a structure with distinguishable element identities.  Element
% colours are neutral (grey / blue / sky / purple) so they cannot be read
% as the semantic green (value) or orange (teacher).
\newcommand{\fMolFinal}{%
  \fBond{ 0.000}{ 0.30}{-0.285}{ 0.09}\fBond{-0.285}{ 0.09}{-0.176}{-0.25}%
  \fBond{-0.176}{-0.25}{ 0.176}{-0.25}\fBond{ 0.176}{-0.25}{ 0.285}{ 0.09}%
  \fBond{ 0.285}{ 0.09}{ 0.000}{ 0.30}%
  \fBond{ 0.285}{ 0.09}{ 0.600}{ 0.27}%
  \fBond{-0.285}{ 0.09}{-0.560}{ 0.32}%
  \fBond{-0.176}{-0.25}{-0.385}{-0.47}%
  \fAtom{ 0.000}{ 0.30}{oiGreyD!75}\fAtom{-0.285}{ 0.09}{oiGreyD!75}%
  \fAtom{-0.176}{-0.25}{oiGreyD!75}\fAtom{ 0.176}{-0.25}{oiGreyD!75}%
  \fAtom{ 0.285}{ 0.09}{oiBlue}%
  \fAtom{ 0.600}{ 0.27}{oiPurple}%
  \fAtomS{-0.560}{ 0.32}{oiSky}%
  \fAtomS{-0.385}{-0.47}{white}}

% the same structure, small, for the displacement cloud in panel (b)
\newcommand{\fMolSmall}[1]{%
  \draw[#1, line width=0.55pt]
    ( 0.000, 0.24) -- (-0.228, 0.07) -- (-0.141,-0.20) --
    ( 0.141,-0.20) -- ( 0.228, 0.07) -- cycle;
  \draw[#1, line width=0.55pt] ( 0.228, 0.07) -- ( 0.470, 0.22);
  \draw[#1, line width=0.55pt] (-0.228, 0.07) -- (-0.440, 0.25);
  \foreach \px/\py in {0/0.24, -0.228/0.07, -0.141/-0.20, 0.141/-0.20,
                       0.228/0.07, 0.470/0.22, -0.440/0.25}
    {\fill[#1] (\px,\py) circle (1.15pt);}}

% =====================================================================
\begin{tikzpicture}[x=1cm, y=1cm, scale=\figOneScale, >=stealth']
% stable bounding box
\path (0,0) rectangle (13.90,8.80);

% ---- panel frames ----------------------------------------------------
\draw[fpanel=oiBlue]   (0.02,4.85) rectangle (6.82,8.75);
\draw[fpanel=oiOrange] (7.10,4.85) rectangle (13.88,8.75);
\draw[fpanel=oiGreen]  (0.02,0.05) rectangle (6.82,4.55);
\draw[fpanel=oiGrey]   (7.10,0.05) rectangle (13.88,4.55);

% =====================================================================
% (a) THE GENERATOR
% =====================================================================
\node[ftitle] at (0.22,8.66) {(a) A de novo flow};
\node[fsub, text width=6.30cm] at (0.22,8.30)
     {composition \emph{and} geometry from one model};

% prior
\begin{scope}[shift={(1.15,7.34)}]\fMolNoise\end{scope}
\node[fctr, text width=2.60cm] at (1.32,6.66)
     {$x_0\sim\mathcal{N}(0,\sigma^{2}I)$};
\node[fctrn, text width=2.40cm] at (1.32,6.28) {random types};

% flow arrow
\draw[flow] (2.20,7.34) -- (3.90,7.34);
\node[anchor=south, font=\footnotesize, text=oiBlue!80!black] at (3.05,7.46)
     {$v_\theta(x_t,t)$};
\node[anchor=north, font=\scriptsize, text=oiGreyD] at (3.05,7.20)
     {flow matching};

% sample
\begin{scope}[shift={(5.15,7.34)}]\fMolFinal\end{scope}
\node[fctr, text width=2.60cm] at (5.15,6.66) {$(c,\,x)$};
\node[fctrn, text width=3.10cm] at (5.15,6.30)
     {atom types $c$, coordinates $x$};

\node[fbody, text width=6.46cm] at (0.22,5.94)
     {No bond labels anywhere. With the discrete labels
      clamped, reversing the positional ODE gives
      $\log q_\theta(x\mid c)$.};

% =====================================================================
% (b) THE TEACHER
% =====================================================================
\node[ftitle] at (7.30,8.66) {(b) Teacher energy values};
\node[fsub, text width=6.30cm] at (7.30,8.30)
     {$K$ displaced geometries of \emph{one} parent};

% displacement cloud: two faint copies + one dark reference structure
\draw[oiGrey!45, line width=0.5pt, dash pattern=on 1.8pt off 1.6pt]
      (8.35,7.24) circle (0.72);
\begin{scope}[shift={(8.19,7.40)}]\fMolSmall{oiGreyL}\end{scope}
\begin{scope}[shift={(8.54,7.10)}]\fMolSmall{oiGreyL}\end{scope}
\begin{scope}[shift={(8.33,7.22)}]\fMolSmall{oiGreyD!80}\end{scope}
\node[fctr, text width=2.40cm] at (8.35,6.36) {$x_1,\dots,x_K$};
\node[fctrn, text width=2.60cm] at (8.35,6.04) {same composition $c$};

% teacher arrow
\draw[teach] (9.28,7.24) -- (10.22,7.24);
\node[anchor=south, font=\footnotesize, text=oiOrange!55!black] at (9.75,7.36)
     {$E(c,\cdot)$};

% energy ladder
\draw[-stealth', line width=0.7pt, draw=oiGreyD] (10.72,6.42) -- (10.72,7.94);
\node[anchor=east, font=\footnotesize, text=oiGreyD] at (10.62,7.84) {$E$};
\foreach \yy in {6.58, 6.86, 7.12, 7.44, 7.74}
  {\draw[oiGrey!70, line width=0.5pt] (10.72,\yy) -- (11.06,\yy);
   \fill[oiOrange!85!black] (11.06,\yy) circle (1.6pt);}
\node[anchor=west, font=\footnotesize] at (11.24,7.62)
     {$E_k=E(c,x_k)$};
\node[anchor=north west, font=\scriptsize, text width=2.50cm, text=oiGreyD]
      at (11.24,7.30) {one value per geometry};

\node[fbody, text width=6.34cm] at (7.30,5.62)
     {A universal neural potential over $83$ elements supplies every
      value. Only the values are used.};

% =====================================================================
% (c) THE OBJECTIVE
% =====================================================================
\node[ftitle] at (0.22,4.46) {(c) Proposed energy-value supervision};

\draw[eqbox] (0.16,2.86) rectangle (6.68,4.06);
\node[anchor=center, font=\small] at (3.42,3.86)
     {$\mathcal{L}_{\mathrm{ours}}
       =\mathcal{L}_{\mathrm{FM}}+\lambda_E\,\mathcal{L}_{\mathrm{value}}$};
\node[anchor=center, font=\small] at (3.42,3.48)
     {$\mathcal{L}_{\mathrm{value}}
       =\mathrm{Var}_{k}\big[\,\log q_\theta(x_k\mid c)+E_k/kT\,\big]$};
% The cancellation is the selling point of the objective, so it is drawn in
% the value colour rather than in neutral grey.  It is a statement about the
% objective, not a measured claim.
\node[anchor=north, font=\scriptsize, inner sep=0pt, align=center,
      text=oiGreen!45!black, text width=6.20cm] at (3.42,3.22)
     {one added term; the partition function cancels};

% schematic alignment plot
\draw[oiGreyD, line width=0.7pt] (1.36,2.52) -- (1.36,1.02) -- (3.98,1.02);
\node[rotate=90, anchor=south, font=\footnotesize, text=oiGreyD]
      at (1.20,1.77) {$\log q_\theta$};
\node[anchor=north, font=\footnotesize, text=oiGreyD] at (2.70,0.94)
     {$-E_k/kT$};
% the line the value term asks for: slope one, offset free
\draw[oiGreen!70!black, line width=0.9pt, dash pattern=on 2.6pt off 1.8pt]
      (1.52,1.16) -- (3.82,2.40);
% unconstrained (grey) and aligned (green) geometries
\foreach \px/\py in {1.68/2.12, 2.02/1.24, 2.36/2.32, 2.72/1.42,
                     3.08/1.98, 3.46/1.30, 3.72/2.20}
   {\fill[oiGrey!75] (\px,\py) circle (1.7pt);}
\foreach \px/\py in {1.68/1.22, 2.02/1.42, 2.36/1.50, 2.72/1.76,
                     3.08/1.92, 3.46/2.14, 3.72/2.36}
   {\fill[oiGreen] (\px,\py) circle (1.7pt);}
% legend
\fill[oiGrey!75] (4.42,2.32) circle (1.7pt);
\node[anchor=west, font=\footnotesize, text width=2.00cm] at (4.58,2.32)
     {$\mathcal{L}_{\mathrm{value}}$ large};
\fill[oiGreen] (4.42,1.86) circle (1.7pt);
\node[anchor=west, font=\footnotesize, text width=2.00cm] at (4.58,1.86)
     {$\mathcal{L}_{\mathrm{value}}=0$};
\node[fnote, text width=2.30cm] at (4.30,1.50)
     {one variance per group};

\node[anchor=north, font=\scriptsize, text=oiGreyD] at (3.42,0.56)
     {$\mathcal{L}_{\mathrm{value}}=0$ within a group
      $\iff \log q_\theta=-E_k/kT+b$};

% =====================================================================
% (d) THE SCOPE
% =====================================================================
% Two lines of \small\bfseries; the composition fact that used to sit in a
% separate subtitle is folded into the region label below, which is where the
% single conditional density is named anyway.
\node[ftitle, text width=6.50cm, align=left] at (7.30,4.46)
     {(d) Local supervision and the\\path to global control};

% the enclosing region: ONE conditional density, several supervised groups
\draw[rounded corners=2.8pt, draw=oiGrey!50, line width=0.6pt,
      dash pattern=on 2.2pt off 1.8pt] (7.42,1.02) rectangle (13.56,3.54);
\node[anchor=center, font=\scriptsize, text=oiGreyD, inner sep=1.6pt,
      fill=oiGrey!4] at (10.49,3.54)
     {one density $q_\theta(\cdot\mid c)$ --- same composition $c$};

\foreach \cx/\lbl in {8.86/A, 12.12/B}
{%
  \node[anchor=north, font=\footnotesize] at (\cx,3.36)
       {group $\lbl$ (basin $\lbl$)};
  \draw[oiGreen!60!white, line width=0.8pt, fill=oiGreen!7]
        (\cx,2.38) ellipse (1.28 and 0.56);
  \draw[oiGreen!45!black, line width=0.6pt]
        (\cx-0.88,2.26) -- (\cx-0.44,2.52) -- (\cx,2.36)
        -- (\cx+0.44,2.69) -- (\cx+0.84,2.50);
  \foreach \dx/\yy in {-0.88/2.26, -0.44/2.52, 0.00/2.36,
                        0.44/2.69, 0.84/2.50}
    {\fill[oiGreen] (\cx+\dx,\yy) circle (1.7pt);}
  \node[anchor=north, font=\scriptsize, text=oiGreyD] at (\cx,1.74)
       {local relation constrained};
  \node[anchor=north, font=\footnotesize] at (\cx,1.46)
       {offset $b_{\lbl}$ free};
}

\draw[freearr] (10.16,2.38) -- (10.82,2.38);
\node[anchor=south, font=\footnotesize, text=oiGrey] at (10.49,2.46) {?};

% Same mathematical content as before -- b_B - b_A and the relative mass are
% not identified -- with the constructive half of the statement added: what
% would have to be built to identify them.  No claim is weakened.
% WIDTH BUDGET.  The float sets \figOneScale = 1.000 and the picture is measured
% at that scale (see the measurement note in fig1_objective_float.tex: at 1.000 a
% 600 dpi render has zero ink outside the four panel frames; at 0.86 it has 4085
% such pixels, because tikz's `scale' shrinks COORDINATES but not type).  Panel
% (d) is 6.78cm wide at scale 1.000 and every font size here is absolute.  Two lines of at most ~42 \footnotesize characters is the whole
% budget; anything longer runs out of the panel and into panel (c).  Both qualifications survive the compression: the offset
% DIFFERENCE and the relative mass are the two unidentified quantities, and the
% second line names the extension that would identify them.
\node[anchor=north, font=\footnotesize, align=center]
      at (10.49,0.92)
     {$b_B-b_A$, relative mass: not identified.\\
      Connecting groups would share one offset.};

\end{tikzpicture}

\figOneFinish

  \caption{\textbf{Energy-value supervision without a partition function.}
    Ideal-objective schematic; no measured quantity appears here. The archived implementation does not realise this density construction (\Cref{sec:method:impl}).
    \textbf{(a)}~With the endpoint discrete labels clamped, reversing the
    positional ODE defines the \emph{clamped positional-flow density}
    $\log q_\theta(x \mid c)$, which need not equal the joint generator's
    conditional density. \textbf{(b)}~A universal neural potential values each
    of $K$ displaced geometries of one parent. \textbf{(c)}~The proposed
    objective is flow matching plus the within-group variance of
    $\log q_\theta(x_k \mid c) + E_k/kT$, which vanishes exactly when
    $\log q_\theta = -E_k/kT + b$ inside the group: the partition function
    cancels and the offset $b$ is free. \textbf{(d)}~Two disconnected groups of
    the \emph{same} density keep independent offsets, so the term constrains
    the local density--energy relation inside each group and leaves their
    relative mass unidentified; connecting groups so they share one offset is
    the extension toward global control. Panel detail in \Cref{app:figdetail}.}
  \label{fig:objective}
\end{figure}

%       or, if the writer prefers to own the float:
%         \begin{figure}[t]\centering
%           % =====================================================================
%  fig1_objective.tex --- Figure 1: the objective and its scope
%
%  CONCEPTUAL FIGURE.  NO DATA.  No number in this figure is measured;
%  every glyph is a schematic.  Nothing here asserts a result.
%
%  Replaces the old five-band figures/fig1_method.tex, which was too
%  dense for a first-page explanatory figure and whose bands (b)/(c)
%  gave the gradient-supervision path equal visual status.  Per the
%  build brief the force branch is REMOVED from this figure entirely;
%  gradient supervision appears only later, as an experimental
%  comparator, and is not drawn here.
%
%  Pure TikZ.  No external images, no pgfplots, no \includegraphics,
%  no arrows.meta, no pgfmathrandom (byte-reproducible across compiles).
%  Okabe--Ito colour-blind-safe palette, semantics shared with
%  make_experiment_figures.py::PALETTE:
%     blue    model / flow / learned quantities
%     green   the value term (energy values, L_value)
%     orange  the teacher (universal neural potential)
%     grey    neutral structure
%  Vermilion (= gradient supervision) and purple (= scrambled control)
%  are deliberately ABSENT: neither concept appears in this figure.
%
%  USAGE
%   (a) standalone (verification / slides):   pdflatex fig1_objective.tex
%   (b) inside the paper -- main.tex already \usepackage{tikz} and
%       \usepackage{amsmath}:
%         % =====================================================================
%  fig1_objective_float.tex --- the float wrapper for Figure 1.
%
%  Figure 1 is pure TikZ, so this wrapper \input's the picture directly
%  rather than \includegraphics-ing a PDF (the same mechanism the old
%  fig1_method.tex used).  main.tex already loads tikz and amsmath.
%
%  Include from sections/01_intro.tex with:
%      % =====================================================================
%  fig1_objective_float.tex --- the float wrapper for Figure 1.
%
%  Figure 1 is pure TikZ, so this wrapper \input's the picture directly
%  rather than \includegraphics-ing a PDF (the same mechanism the old
%  fig1_method.tex used).  main.tex already loads tikz and amsmath.
%
%  Include from sections/01_intro.tex with:
%      \input{figures/fig1_objective_float}
%
%  The picture is drawn at 13.90cm = 5.47in, i.e. full ICLR single-column
%  width at scale 1.0.  Do NOT wrap it in \resizebox and do NOT scale it:
%  the in-figure type is \small / \footnotesize / \scriptsize of the 10pt
%  body font and scaling would desynchronise it from the caption.
%
%  PROVENANCE.  Conceptual figure, no data.  Built per the Figure 1 brief:
%  four panels (generator / teacher values / the value term / scope), and
%  the gradient-supervision branch of the superseded fig1_method.tex is
%  removed entirely rather than demoted.  No panel asserts a result.
%
%  REVISION C13 (figure-editor pass).  The framing moves to the method: panel
%  (c) is titled "Proposed energy-value supervision" and states inside the box
%  that the partition function cancels; panel (d)'s heading is "Local
%  supervision and the path to global control" and its closing line names the
%  extension (connecting groups so they share one offset) alongside the
%  unchanged qualification that b_B - b_A and the relative mass are not
%  identified.  No mathematical content was removed.
%
%  REVISION C8.  The caption is rewritten only where panels (c) and (d)
%  changed.  (c) now names the complete objective, so the caption states
%  it.  (d) now shows two disconnected groups of ONE composition, so the
%  caption no longer says the offset is free "between compositions"; the
%  freedom is already there inside a single conditional density.  The
%  verb "pins" is gone: the term CONSTRAINS THE LOCAL RELATION between
%  log-density and energy within a group.  The disconnection claim is
%  empirical and audited -- distinct parents sharing a composition supply
%  distinct perturbation clouds that share no sampled geometry (31% of
%  training parents sit on a repeated composition and are still
%  disconnected), so the overlap graph is a set of isolated vertices.
% =====================================================================
\begin{figure}[t]
  \centering
  \providecommand{\figOneScale}{1.0}\renewcommand{\figOneScale}{1.000}%
% SCALE = 1.000, THE FIGURE's NATURAL SIZE.  Do not lower it to buy page space.
% TikZ's `scale' multiplies COORDINATES but not type: every label in this
% picture is set at an absolute size (\small / \footnotesize / \scriptsize of
% the 10pt body), so any scale < 1 shrinks the panel frames out from under
% text that keeps its size, and the text runs out of the panels.
% MEASURED 2026-08-30 (600 dpi render of the standalone picture; panel frames
% located from the rendered strokes, so the calibration is self-checking at
% 236.22 px/cm on both axes):
%     scale 0.86 -> 4085 ink pixels outside the four panel frames, in 6
%                   regions.  Panel (b)'s body line runs to x = 14.59 cm,
%                   i.e. 0.71 cm past its own frame and 0.69 cm past the
%                   picture's 13.90 cm bounding box; panel (a)'s last two lines
%                   cross into panel (b)'s frame; panel (c)'s closing identity
%                   line crosses the left frame at x = -0.17 cm.
%     scale 1.000 -> 0 ink pixels outside the panels, at every ink threshold
%                   tested (lum < 128 and lum < 200, with 5 px of stroke
%                   tolerance).  All four panels are clean.
% PAGE COST: none that changes the budget.  With 1.000 the main text ends on
% page 9 and a probe \vspace inserted before \label{endofmaintext} shows a
% further +108pt fits before it spills to page 10 (+114pt spills); at 0.86 the
% same probe breaks between +108pt and +120pt.  The fix costs at most about
% one body line of final-page headroom.
  \input{figures/fig1_objective}
  \caption{\textbf{Energy-value supervision without a partition function.}
    Ideal-objective schematic; no measured quantity appears here. The archived implementation does not realise this density construction (\Cref{sec:method:impl}).
    \textbf{(a)}~With the endpoint discrete labels clamped, reversing the
    positional ODE defines the \emph{clamped positional-flow density}
    $\log q_\theta(x \mid c)$, which need not equal the joint generator's
    conditional density. \textbf{(b)}~A universal neural potential values each
    of $K$ displaced geometries of one parent. \textbf{(c)}~The proposed
    objective is flow matching plus the within-group variance of
    $\log q_\theta(x_k \mid c) + E_k/kT$, which vanishes exactly when
    $\log q_\theta = -E_k/kT + b$ inside the group: the partition function
    cancels and the offset $b$ is free. \textbf{(d)}~Two disconnected groups of
    the \emph{same} density keep independent offsets, so the term constrains
    the local density--energy relation inside each group and leaves their
    relative mass unidentified; connecting groups so they share one offset is
    the extension toward global control. Panel detail in \Cref{app:figdetail}.}
  \label{fig:objective}
\end{figure}

%
%  The picture is drawn at 13.90cm = 5.47in, i.e. full ICLR single-column
%  width at scale 1.0.  Do NOT wrap it in \resizebox and do NOT scale it:
%  the in-figure type is \small / \footnotesize / \scriptsize of the 10pt
%  body font and scaling would desynchronise it from the caption.
%
%  PROVENANCE.  Conceptual figure, no data.  Built per the Figure 1 brief:
%  four panels (generator / teacher values / the value term / scope), and
%  the gradient-supervision branch of the superseded fig1_method.tex is
%  removed entirely rather than demoted.  No panel asserts a result.
%
%  REVISION C13 (figure-editor pass).  The framing moves to the method: panel
%  (c) is titled "Proposed energy-value supervision" and states inside the box
%  that the partition function cancels; panel (d)'s heading is "Local
%  supervision and the path to global control" and its closing line names the
%  extension (connecting groups so they share one offset) alongside the
%  unchanged qualification that b_B - b_A and the relative mass are not
%  identified.  No mathematical content was removed.
%
%  REVISION C8.  The caption is rewritten only where panels (c) and (d)
%  changed.  (c) now names the complete objective, so the caption states
%  it.  (d) now shows two disconnected groups of ONE composition, so the
%  caption no longer says the offset is free "between compositions"; the
%  freedom is already there inside a single conditional density.  The
%  verb "pins" is gone: the term CONSTRAINS THE LOCAL RELATION between
%  log-density and energy within a group.  The disconnection claim is
%  empirical and audited -- distinct parents sharing a composition supply
%  distinct perturbation clouds that share no sampled geometry (31% of
%  training parents sit on a repeated composition and are still
%  disconnected), so the overlap graph is a set of isolated vertices.
% =====================================================================
\begin{figure}[t]
  \centering
  \providecommand{\figOneScale}{1.0}\renewcommand{\figOneScale}{1.000}%
% SCALE = 1.000, THE FIGURE's NATURAL SIZE.  Do not lower it to buy page space.
% TikZ's `scale' multiplies COORDINATES but not type: every label in this
% picture is set at an absolute size (\small / \footnotesize / \scriptsize of
% the 10pt body), so any scale < 1 shrinks the panel frames out from under
% text that keeps its size, and the text runs out of the panels.
% MEASURED 2026-08-30 (600 dpi render of the standalone picture; panel frames
% located from the rendered strokes, so the calibration is self-checking at
% 236.22 px/cm on both axes):
%     scale 0.86 -> 4085 ink pixels outside the four panel frames, in 6
%                   regions.  Panel (b)'s body line runs to x = 14.59 cm,
%                   i.e. 0.71 cm past its own frame and 0.69 cm past the
%                   picture's 13.90 cm bounding box; panel (a)'s last two lines
%                   cross into panel (b)'s frame; panel (c)'s closing identity
%                   line crosses the left frame at x = -0.17 cm.
%     scale 1.000 -> 0 ink pixels outside the panels, at every ink threshold
%                   tested (lum < 128 and lum < 200, with 5 px of stroke
%                   tolerance).  All four panels are clean.
% PAGE COST: none that changes the budget.  With 1.000 the main text ends on
% page 9 and a probe \vspace inserted before \label{endofmaintext} shows a
% further +108pt fits before it spills to page 10 (+114pt spills); at 0.86 the
% same probe breaks between +108pt and +120pt.  The fix costs at most about
% one body line of final-page headroom.
  % =====================================================================
%  fig1_objective.tex --- Figure 1: the objective and its scope
%
%  CONCEPTUAL FIGURE.  NO DATA.  No number in this figure is measured;
%  every glyph is a schematic.  Nothing here asserts a result.
%
%  Replaces the old five-band figures/fig1_method.tex, which was too
%  dense for a first-page explanatory figure and whose bands (b)/(c)
%  gave the gradient-supervision path equal visual status.  Per the
%  build brief the force branch is REMOVED from this figure entirely;
%  gradient supervision appears only later, as an experimental
%  comparator, and is not drawn here.
%
%  Pure TikZ.  No external images, no pgfplots, no \includegraphics,
%  no arrows.meta, no pgfmathrandom (byte-reproducible across compiles).
%  Okabe--Ito colour-blind-safe palette, semantics shared with
%  make_experiment_figures.py::PALETTE:
%     blue    model / flow / learned quantities
%     green   the value term (energy values, L_value)
%     orange  the teacher (universal neural potential)
%     grey    neutral structure
%  Vermilion (= gradient supervision) and purple (= scrambled control)
%  are deliberately ABSENT: neither concept appears in this figure.
%
%  USAGE
%   (a) standalone (verification / slides):   pdflatex fig1_objective.tex
%   (b) inside the paper -- main.tex already \usepackage{tikz} and
%       \usepackage{amsmath}:
%         \input{figures/fig1_objective_float}
%       or, if the writer prefers to own the float:
%         \begin{figure}[t]\centering
%           \input{figures/fig1_objective}
%           \caption{...}\label{fig:objective}
%         \end{figure}
%
%  SIZE.  Natural size 13.90cm x 8.80cm == 5.47in x 3.46in.
%  ICLR is a SINGLE 5.5in text column, so this is a full-text-width
%  figure at scale 1.0.  DO NOT rescale it and DO NOT wrap it in
%  \resizebox: every glyph inside is \small / \footnotesize /
%  \scriptsize of the 10pt body font, and scaling desynchronises
%  in-figure text from the caption.
%
%  LAYOUT.  A 2x2 grid, read left-to-right, top-to-bottom:
%    (a) the generator      composition AND geometry from one flow
%    (b) the teacher        energy values for K displacements of one parent
%    (c) the objective      L_ours = L_FM + lambda_E L_value, and nothing else
%    (d) the scope          two basins of ONE composition; offset free
%  Panel widths 6.80 / 6.78 cm; no wire crosses a panel boundary, so
%  there is no routing to collide with text.
%
%  REVISION C8 (this pass).  Two changes, both confined to panels (c)/(d):
%   * (c) now displays the complete objective, L_ours = L_FM + lambda_E
%     L_value, so the reader cannot infer a hidden force or anchor branch.
%     Neither branch was ever drawn here and neither is added.
%   * (d) previously drew two groups of two DIFFERENT compositions
%     (c_A, c_B).  That was wrong twice over.  It misstated the failure
%     mode -- the unidentified offset is already present WITHIN a single
%     conditional density -- and it flirted with the claim that one
%     composition contributes exactly one group, which the composition
%     audit refutes: in the training perturbation shard
%     (processed_data/omol25_4m_processed/perturbation_train_n30000_s0.pt)
%     30,000 parent groups carry only
%     22,822 distinct (element multiset, total charge) keys, 2,204 keys
%     are shared by more than one parent (9,382 parents, 31%; max
%     multiplicity 80).  Panel (d) therefore now shows two groups A and B
%     of the SAME composition c: two disconnected supervised regions of
%     one density p_theta(.|c), each with its own free offset b_A, b_B.
%     "ordering fixed" / "pinned" wording is replaced throughout by
%     "the local relation between log-density and energy is constrained".
%  In the EVALUATION population (val_data_processed.pt; the 120 grid
%  parents and the 93 primary parents) every parent has a distinct
%  composition, multiplicity 1 -- that is an audited property of the
%  evaluation set, not of the supervision design, and this figure does
%  not assert it.
% =====================================================================

% ---- standalone bootstrap (skipped when \input into a tikz-enabled doc) ----
\ifx\tikzpicture\undefined
  \documentclass[10pt]{article}
  % 0.05cm of slack past the 13.90cm bounding box so that the panel border
  % stroke width does not report a spurious overfull \hbox in standalone mode.
  \usepackage[paperwidth=14.40cm,paperheight=9.20cm,margin=0.20cm,noheadfoot]{geometry}
  \usepackage{amsmath,amssymb,bm}
  \usepackage{times}
  \usepackage{tikz}
  \setlength{\topskip}{0pt}
  \setlength{\parindent}{0pt}
  \pagestyle{empty}
  \begin{document}%
  \def\figOneFinish{\end{document}}
\fi
\providecommand{\figOneFinish}{\relax}
\providecommand{\figOneScale}{1.0}

\usetikzlibrary{positioning,calc,fit,backgrounds,arrows,shapes.geometric}

% ---------------------------------------------------------------- palette
% Okabe & Ito (2008), CVD-safe qualitative set.
\providecommand{\oiOrangeDefined}{}
\definecolor{oiOrange}{HTML}{E69F00}
\definecolor{oiSky}   {HTML}{56B4E9}
\definecolor{oiGreen} {HTML}{009E73}
\definecolor{oiBlue}  {HTML}{0072B2}
\definecolor{oiPurple}{HTML}{CC79A7}
\definecolor{oiGrey}  {HTML}{7F7F7F}
\definecolor{oiGreyL} {HTML}{C4C4C4}
\definecolor{oiGreyD} {HTML}{4D4D4D}

% ---------------------------------------------------------------- type
% Body 10pt Times => \small 9pt, \footnotesize 8pt, \scriptsize 7pt.
% Titles are 9pt bold; all explanatory text is 8pt; 7pt is used only for
% short qualifiers.  Nothing in this figure is smaller than 7pt.
\tikzset{
  fpanel/.style   ={rounded corners=2.4pt, draw=#1!55!white, line width=0.7pt,
                    fill=#1!4},
  ftitle/.style   ={anchor=north west, font=\small\bfseries, inner sep=0pt,
                    text=oiGreyD},
  fsub/.style     ={anchor=north west, font=\footnotesize, inner sep=0pt,
                    text=oiGrey, align=left},
  fbody/.style    ={anchor=north west, font=\footnotesize, inner sep=0pt,
                    align=left},
  fnote/.style    ={anchor=north west, font=\scriptsize, inner sep=0pt,
                    align=left, text=oiGreyD},
  fctr/.style     ={anchor=north, font=\footnotesize, inner sep=0pt,
                    align=center},
  fctrn/.style    ={anchor=north, font=\scriptsize, inner sep=0pt,
                    align=center, text=oiGreyD},
  flow/.style     ={-stealth', line width=1.0pt, draw=oiBlue},
  teach/.style    ={-stealth', line width=1.0pt, draw=oiOrange!88!black},
  freearr/.style  ={stealth'-stealth', line width=0.8pt, draw=oiGrey,
                    dash pattern=on 2.4pt off 1.8pt},
  eqbox/.style    ={rounded corners=2.4pt, draw=oiGreen!65!white,
                    line width=0.8pt, fill=oiGreen!8},
}

% ------------------------------------------------------- molecule glyphs
% Deterministic coordinates: the figure is identical on every compile.
\newcommand{\fAtom}[3]{\fill[#3] (#1,#2) circle (2.05pt);%
  \draw[oiGreyD!55, line width=0.2pt] (#1,#2) circle (2.05pt);}
\newcommand{\fAtomS}[3]{\fill[#3] (#1,#2) circle (1.35pt);%
  \draw[oiGreyD!45, line width=0.15pt] (#1,#2) circle (1.35pt);}
\newcommand{\fBond}[4]{\draw[oiGrey!80, line width=0.7pt] (#1,#2) -- (#3,#4);}

% t = 0 : unstructured cloud, no bonds, no identities
\newcommand{\fMolNoise}{%
  \fAtomS{-0.36}{ 0.32}{oiGreyL}\fAtomS{ 0.12}{ 0.38}{oiGreyL}%
  \fAtomS{ 0.38}{ 0.06}{oiGreyL}\fAtomS{-0.10}{-0.02}{oiGreyL}%
  \fAtomS{-0.42}{-0.28}{oiGreyL}\fAtomS{ 0.06}{-0.36}{oiGreyL}%
  \fAtomS{ 0.42}{-0.34}{oiGreyL}\fAtomS{ 0.26}{ 0.24}{oiGreyL}%
  \fAtomS{-0.18}{ 0.12}{oiGreyL}}

% t = 1 : a structure with distinguishable element identities.  Element
% colours are neutral (grey / blue / sky / purple) so they cannot be read
% as the semantic green (value) or orange (teacher).
\newcommand{\fMolFinal}{%
  \fBond{ 0.000}{ 0.30}{-0.285}{ 0.09}\fBond{-0.285}{ 0.09}{-0.176}{-0.25}%
  \fBond{-0.176}{-0.25}{ 0.176}{-0.25}\fBond{ 0.176}{-0.25}{ 0.285}{ 0.09}%
  \fBond{ 0.285}{ 0.09}{ 0.000}{ 0.30}%
  \fBond{ 0.285}{ 0.09}{ 0.600}{ 0.27}%
  \fBond{-0.285}{ 0.09}{-0.560}{ 0.32}%
  \fBond{-0.176}{-0.25}{-0.385}{-0.47}%
  \fAtom{ 0.000}{ 0.30}{oiGreyD!75}\fAtom{-0.285}{ 0.09}{oiGreyD!75}%
  \fAtom{-0.176}{-0.25}{oiGreyD!75}\fAtom{ 0.176}{-0.25}{oiGreyD!75}%
  \fAtom{ 0.285}{ 0.09}{oiBlue}%
  \fAtom{ 0.600}{ 0.27}{oiPurple}%
  \fAtomS{-0.560}{ 0.32}{oiSky}%
  \fAtomS{-0.385}{-0.47}{white}}

% the same structure, small, for the displacement cloud in panel (b)
\newcommand{\fMolSmall}[1]{%
  \draw[#1, line width=0.55pt]
    ( 0.000, 0.24) -- (-0.228, 0.07) -- (-0.141,-0.20) --
    ( 0.141,-0.20) -- ( 0.228, 0.07) -- cycle;
  \draw[#1, line width=0.55pt] ( 0.228, 0.07) -- ( 0.470, 0.22);
  \draw[#1, line width=0.55pt] (-0.228, 0.07) -- (-0.440, 0.25);
  \foreach \px/\py in {0/0.24, -0.228/0.07, -0.141/-0.20, 0.141/-0.20,
                       0.228/0.07, 0.470/0.22, -0.440/0.25}
    {\fill[#1] (\px,\py) circle (1.15pt);}}

% =====================================================================
\begin{tikzpicture}[x=1cm, y=1cm, scale=\figOneScale, >=stealth']
% stable bounding box
\path (0,0) rectangle (13.90,8.80);

% ---- panel frames ----------------------------------------------------
\draw[fpanel=oiBlue]   (0.02,4.85) rectangle (6.82,8.75);
\draw[fpanel=oiOrange] (7.10,4.85) rectangle (13.88,8.75);
\draw[fpanel=oiGreen]  (0.02,0.05) rectangle (6.82,4.55);
\draw[fpanel=oiGrey]   (7.10,0.05) rectangle (13.88,4.55);

% =====================================================================
% (a) THE GENERATOR
% =====================================================================
\node[ftitle] at (0.22,8.66) {(a) A de novo flow};
\node[fsub, text width=6.30cm] at (0.22,8.30)
     {composition \emph{and} geometry from one model};

% prior
\begin{scope}[shift={(1.15,7.34)}]\fMolNoise\end{scope}
\node[fctr, text width=2.60cm] at (1.32,6.66)
     {$x_0\sim\mathcal{N}(0,\sigma^{2}I)$};
\node[fctrn, text width=2.40cm] at (1.32,6.28) {random types};

% flow arrow
\draw[flow] (2.20,7.34) -- (3.90,7.34);
\node[anchor=south, font=\footnotesize, text=oiBlue!80!black] at (3.05,7.46)
     {$v_\theta(x_t,t)$};
\node[anchor=north, font=\scriptsize, text=oiGreyD] at (3.05,7.20)
     {flow matching};

% sample
\begin{scope}[shift={(5.15,7.34)}]\fMolFinal\end{scope}
\node[fctr, text width=2.60cm] at (5.15,6.66) {$(c,\,x)$};
\node[fctrn, text width=3.10cm] at (5.15,6.30)
     {atom types $c$, coordinates $x$};

\node[fbody, text width=6.46cm] at (0.22,5.94)
     {No bond labels anywhere. With the discrete labels
      clamped, reversing the positional ODE gives
      $\log q_\theta(x\mid c)$.};

% =====================================================================
% (b) THE TEACHER
% =====================================================================
\node[ftitle] at (7.30,8.66) {(b) Teacher energy values};
\node[fsub, text width=6.30cm] at (7.30,8.30)
     {$K$ displaced geometries of \emph{one} parent};

% displacement cloud: two faint copies + one dark reference structure
\draw[oiGrey!45, line width=0.5pt, dash pattern=on 1.8pt off 1.6pt]
      (8.35,7.24) circle (0.72);
\begin{scope}[shift={(8.19,7.40)}]\fMolSmall{oiGreyL}\end{scope}
\begin{scope}[shift={(8.54,7.10)}]\fMolSmall{oiGreyL}\end{scope}
\begin{scope}[shift={(8.33,7.22)}]\fMolSmall{oiGreyD!80}\end{scope}
\node[fctr, text width=2.40cm] at (8.35,6.36) {$x_1,\dots,x_K$};
\node[fctrn, text width=2.60cm] at (8.35,6.04) {same composition $c$};

% teacher arrow
\draw[teach] (9.28,7.24) -- (10.22,7.24);
\node[anchor=south, font=\footnotesize, text=oiOrange!55!black] at (9.75,7.36)
     {$E(c,\cdot)$};

% energy ladder
\draw[-stealth', line width=0.7pt, draw=oiGreyD] (10.72,6.42) -- (10.72,7.94);
\node[anchor=east, font=\footnotesize, text=oiGreyD] at (10.62,7.84) {$E$};
\foreach \yy in {6.58, 6.86, 7.12, 7.44, 7.74}
  {\draw[oiGrey!70, line width=0.5pt] (10.72,\yy) -- (11.06,\yy);
   \fill[oiOrange!85!black] (11.06,\yy) circle (1.6pt);}
\node[anchor=west, font=\footnotesize] at (11.24,7.62)
     {$E_k=E(c,x_k)$};
\node[anchor=north west, font=\scriptsize, text width=2.50cm, text=oiGreyD]
      at (11.24,7.30) {one value per geometry};

\node[fbody, text width=6.34cm] at (7.30,5.62)
     {A universal neural potential over $83$ elements supplies every
      value. Only the values are used.};

% =====================================================================
% (c) THE OBJECTIVE
% =====================================================================
\node[ftitle] at (0.22,4.46) {(c) Proposed energy-value supervision};

\draw[eqbox] (0.16,2.86) rectangle (6.68,4.06);
\node[anchor=center, font=\small] at (3.42,3.86)
     {$\mathcal{L}_{\mathrm{ours}}
       =\mathcal{L}_{\mathrm{FM}}+\lambda_E\,\mathcal{L}_{\mathrm{value}}$};
\node[anchor=center, font=\small] at (3.42,3.48)
     {$\mathcal{L}_{\mathrm{value}}
       =\mathrm{Var}_{k}\big[\,\log q_\theta(x_k\mid c)+E_k/kT\,\big]$};
% The cancellation is the selling point of the objective, so it is drawn in
% the value colour rather than in neutral grey.  It is a statement about the
% objective, not a measured claim.
\node[anchor=north, font=\scriptsize, inner sep=0pt, align=center,
      text=oiGreen!45!black, text width=6.20cm] at (3.42,3.22)
     {one added term; the partition function cancels};

% schematic alignment plot
\draw[oiGreyD, line width=0.7pt] (1.36,2.52) -- (1.36,1.02) -- (3.98,1.02);
\node[rotate=90, anchor=south, font=\footnotesize, text=oiGreyD]
      at (1.20,1.77) {$\log q_\theta$};
\node[anchor=north, font=\footnotesize, text=oiGreyD] at (2.70,0.94)
     {$-E_k/kT$};
% the line the value term asks for: slope one, offset free
\draw[oiGreen!70!black, line width=0.9pt, dash pattern=on 2.6pt off 1.8pt]
      (1.52,1.16) -- (3.82,2.40);
% unconstrained (grey) and aligned (green) geometries
\foreach \px/\py in {1.68/2.12, 2.02/1.24, 2.36/2.32, 2.72/1.42,
                     3.08/1.98, 3.46/1.30, 3.72/2.20}
   {\fill[oiGrey!75] (\px,\py) circle (1.7pt);}
\foreach \px/\py in {1.68/1.22, 2.02/1.42, 2.36/1.50, 2.72/1.76,
                     3.08/1.92, 3.46/2.14, 3.72/2.36}
   {\fill[oiGreen] (\px,\py) circle (1.7pt);}
% legend
\fill[oiGrey!75] (4.42,2.32) circle (1.7pt);
\node[anchor=west, font=\footnotesize, text width=2.00cm] at (4.58,2.32)
     {$\mathcal{L}_{\mathrm{value}}$ large};
\fill[oiGreen] (4.42,1.86) circle (1.7pt);
\node[anchor=west, font=\footnotesize, text width=2.00cm] at (4.58,1.86)
     {$\mathcal{L}_{\mathrm{value}}=0$};
\node[fnote, text width=2.30cm] at (4.30,1.50)
     {one variance per group};

\node[anchor=north, font=\scriptsize, text=oiGreyD] at (3.42,0.56)
     {$\mathcal{L}_{\mathrm{value}}=0$ within a group
      $\iff \log q_\theta=-E_k/kT+b$};

% =====================================================================
% (d) THE SCOPE
% =====================================================================
% Two lines of \small\bfseries; the composition fact that used to sit in a
% separate subtitle is folded into the region label below, which is where the
% single conditional density is named anyway.
\node[ftitle, text width=6.50cm, align=left] at (7.30,4.46)
     {(d) Local supervision and the\\path to global control};

% the enclosing region: ONE conditional density, several supervised groups
\draw[rounded corners=2.8pt, draw=oiGrey!50, line width=0.6pt,
      dash pattern=on 2.2pt off 1.8pt] (7.42,1.02) rectangle (13.56,3.54);
\node[anchor=center, font=\scriptsize, text=oiGreyD, inner sep=1.6pt,
      fill=oiGrey!4] at (10.49,3.54)
     {one density $q_\theta(\cdot\mid c)$ --- same composition $c$};

\foreach \cx/\lbl in {8.86/A, 12.12/B}
{%
  \node[anchor=north, font=\footnotesize] at (\cx,3.36)
       {group $\lbl$ (basin $\lbl$)};
  \draw[oiGreen!60!white, line width=0.8pt, fill=oiGreen!7]
        (\cx,2.38) ellipse (1.28 and 0.56);
  \draw[oiGreen!45!black, line width=0.6pt]
        (\cx-0.88,2.26) -- (\cx-0.44,2.52) -- (\cx,2.36)
        -- (\cx+0.44,2.69) -- (\cx+0.84,2.50);
  \foreach \dx/\yy in {-0.88/2.26, -0.44/2.52, 0.00/2.36,
                        0.44/2.69, 0.84/2.50}
    {\fill[oiGreen] (\cx+\dx,\yy) circle (1.7pt);}
  \node[anchor=north, font=\scriptsize, text=oiGreyD] at (\cx,1.74)
       {local relation constrained};
  \node[anchor=north, font=\footnotesize] at (\cx,1.46)
       {offset $b_{\lbl}$ free};
}

\draw[freearr] (10.16,2.38) -- (10.82,2.38);
\node[anchor=south, font=\footnotesize, text=oiGrey] at (10.49,2.46) {?};

% Same mathematical content as before -- b_B - b_A and the relative mass are
% not identified -- with the constructive half of the statement added: what
% would have to be built to identify them.  No claim is weakened.
% WIDTH BUDGET.  The float sets \figOneScale = 1.000 and the picture is measured
% at that scale (see the measurement note in fig1_objective_float.tex: at 1.000 a
% 600 dpi render has zero ink outside the four panel frames; at 0.86 it has 4085
% such pixels, because tikz's `scale' shrinks COORDINATES but not type).  Panel
% (d) is 6.78cm wide at scale 1.000 and every font size here is absolute.  Two lines of at most ~42 \footnotesize characters is the whole
% budget; anything longer runs out of the panel and into panel (c).  Both qualifications survive the compression: the offset
% DIFFERENCE and the relative mass are the two unidentified quantities, and the
% second line names the extension that would identify them.
\node[anchor=north, font=\footnotesize, align=center]
      at (10.49,0.92)
     {$b_B-b_A$, relative mass: not identified.\\
      Connecting groups would share one offset.};

\end{tikzpicture}

\figOneFinish

  \caption{\textbf{Energy-value supervision without a partition function.}
    Ideal-objective schematic; no measured quantity appears here. The archived implementation does not realise this density construction (\Cref{sec:method:impl}).
    \textbf{(a)}~With the endpoint discrete labels clamped, reversing the
    positional ODE defines the \emph{clamped positional-flow density}
    $\log q_\theta(x \mid c)$, which need not equal the joint generator's
    conditional density. \textbf{(b)}~A universal neural potential values each
    of $K$ displaced geometries of one parent. \textbf{(c)}~The proposed
    objective is flow matching plus the within-group variance of
    $\log q_\theta(x_k \mid c) + E_k/kT$, which vanishes exactly when
    $\log q_\theta = -E_k/kT + b$ inside the group: the partition function
    cancels and the offset $b$ is free. \textbf{(d)}~Two disconnected groups of
    the \emph{same} density keep independent offsets, so the term constrains
    the local density--energy relation inside each group and leaves their
    relative mass unidentified; connecting groups so they share one offset is
    the extension toward global control. Panel detail in \Cref{app:figdetail}.}
  \label{fig:objective}
\end{figure}

%       or, if the writer prefers to own the float:
%         \begin{figure}[t]\centering
%           % =====================================================================
%  fig1_objective.tex --- Figure 1: the objective and its scope
%
%  CONCEPTUAL FIGURE.  NO DATA.  No number in this figure is measured;
%  every glyph is a schematic.  Nothing here asserts a result.
%
%  Replaces the old five-band figures/fig1_method.tex, which was too
%  dense for a first-page explanatory figure and whose bands (b)/(c)
%  gave the gradient-supervision path equal visual status.  Per the
%  build brief the force branch is REMOVED from this figure entirely;
%  gradient supervision appears only later, as an experimental
%  comparator, and is not drawn here.
%
%  Pure TikZ.  No external images, no pgfplots, no \includegraphics,
%  no arrows.meta, no pgfmathrandom (byte-reproducible across compiles).
%  Okabe--Ito colour-blind-safe palette, semantics shared with
%  make_experiment_figures.py::PALETTE:
%     blue    model / flow / learned quantities
%     green   the value term (energy values, L_value)
%     orange  the teacher (universal neural potential)
%     grey    neutral structure
%  Vermilion (= gradient supervision) and purple (= scrambled control)
%  are deliberately ABSENT: neither concept appears in this figure.
%
%  USAGE
%   (a) standalone (verification / slides):   pdflatex fig1_objective.tex
%   (b) inside the paper -- main.tex already \usepackage{tikz} and
%       \usepackage{amsmath}:
%         % =====================================================================
%  fig1_objective_float.tex --- the float wrapper for Figure 1.
%
%  Figure 1 is pure TikZ, so this wrapper \input's the picture directly
%  rather than \includegraphics-ing a PDF (the same mechanism the old
%  fig1_method.tex used).  main.tex already loads tikz and amsmath.
%
%  Include from sections/01_intro.tex with:
%      \input{figures/fig1_objective_float}
%
%  The picture is drawn at 13.90cm = 5.47in, i.e. full ICLR single-column
%  width at scale 1.0.  Do NOT wrap it in \resizebox and do NOT scale it:
%  the in-figure type is \small / \footnotesize / \scriptsize of the 10pt
%  body font and scaling would desynchronise it from the caption.
%
%  PROVENANCE.  Conceptual figure, no data.  Built per the Figure 1 brief:
%  four panels (generator / teacher values / the value term / scope), and
%  the gradient-supervision branch of the superseded fig1_method.tex is
%  removed entirely rather than demoted.  No panel asserts a result.
%
%  REVISION C13 (figure-editor pass).  The framing moves to the method: panel
%  (c) is titled "Proposed energy-value supervision" and states inside the box
%  that the partition function cancels; panel (d)'s heading is "Local
%  supervision and the path to global control" and its closing line names the
%  extension (connecting groups so they share one offset) alongside the
%  unchanged qualification that b_B - b_A and the relative mass are not
%  identified.  No mathematical content was removed.
%
%  REVISION C8.  The caption is rewritten only where panels (c) and (d)
%  changed.  (c) now names the complete objective, so the caption states
%  it.  (d) now shows two disconnected groups of ONE composition, so the
%  caption no longer says the offset is free "between compositions"; the
%  freedom is already there inside a single conditional density.  The
%  verb "pins" is gone: the term CONSTRAINS THE LOCAL RELATION between
%  log-density and energy within a group.  The disconnection claim is
%  empirical and audited -- distinct parents sharing a composition supply
%  distinct perturbation clouds that share no sampled geometry (31% of
%  training parents sit on a repeated composition and are still
%  disconnected), so the overlap graph is a set of isolated vertices.
% =====================================================================
\begin{figure}[t]
  \centering
  \providecommand{\figOneScale}{1.0}\renewcommand{\figOneScale}{1.000}%
% SCALE = 1.000, THE FIGURE's NATURAL SIZE.  Do not lower it to buy page space.
% TikZ's `scale' multiplies COORDINATES but not type: every label in this
% picture is set at an absolute size (\small / \footnotesize / \scriptsize of
% the 10pt body), so any scale < 1 shrinks the panel frames out from under
% text that keeps its size, and the text runs out of the panels.
% MEASURED 2026-08-30 (600 dpi render of the standalone picture; panel frames
% located from the rendered strokes, so the calibration is self-checking at
% 236.22 px/cm on both axes):
%     scale 0.86 -> 4085 ink pixels outside the four panel frames, in 6
%                   regions.  Panel (b)'s body line runs to x = 14.59 cm,
%                   i.e. 0.71 cm past its own frame and 0.69 cm past the
%                   picture's 13.90 cm bounding box; panel (a)'s last two lines
%                   cross into panel (b)'s frame; panel (c)'s closing identity
%                   line crosses the left frame at x = -0.17 cm.
%     scale 1.000 -> 0 ink pixels outside the panels, at every ink threshold
%                   tested (lum < 128 and lum < 200, with 5 px of stroke
%                   tolerance).  All four panels are clean.
% PAGE COST: none that changes the budget.  With 1.000 the main text ends on
% page 9 and a probe \vspace inserted before \label{endofmaintext} shows a
% further +108pt fits before it spills to page 10 (+114pt spills); at 0.86 the
% same probe breaks between +108pt and +120pt.  The fix costs at most about
% one body line of final-page headroom.
  \input{figures/fig1_objective}
  \caption{\textbf{Energy-value supervision without a partition function.}
    Ideal-objective schematic; no measured quantity appears here. The archived implementation does not realise this density construction (\Cref{sec:method:impl}).
    \textbf{(a)}~With the endpoint discrete labels clamped, reversing the
    positional ODE defines the \emph{clamped positional-flow density}
    $\log q_\theta(x \mid c)$, which need not equal the joint generator's
    conditional density. \textbf{(b)}~A universal neural potential values each
    of $K$ displaced geometries of one parent. \textbf{(c)}~The proposed
    objective is flow matching plus the within-group variance of
    $\log q_\theta(x_k \mid c) + E_k/kT$, which vanishes exactly when
    $\log q_\theta = -E_k/kT + b$ inside the group: the partition function
    cancels and the offset $b$ is free. \textbf{(d)}~Two disconnected groups of
    the \emph{same} density keep independent offsets, so the term constrains
    the local density--energy relation inside each group and leaves their
    relative mass unidentified; connecting groups so they share one offset is
    the extension toward global control. Panel detail in \Cref{app:figdetail}.}
  \label{fig:objective}
\end{figure}

%       or, if the writer prefers to own the float:
%         \begin{figure}[t]\centering
%           % =====================================================================
%  fig1_objective.tex --- Figure 1: the objective and its scope
%
%  CONCEPTUAL FIGURE.  NO DATA.  No number in this figure is measured;
%  every glyph is a schematic.  Nothing here asserts a result.
%
%  Replaces the old five-band figures/fig1_method.tex, which was too
%  dense for a first-page explanatory figure and whose bands (b)/(c)
%  gave the gradient-supervision path equal visual status.  Per the
%  build brief the force branch is REMOVED from this figure entirely;
%  gradient supervision appears only later, as an experimental
%  comparator, and is not drawn here.
%
%  Pure TikZ.  No external images, no pgfplots, no \includegraphics,
%  no arrows.meta, no pgfmathrandom (byte-reproducible across compiles).
%  Okabe--Ito colour-blind-safe palette, semantics shared with
%  make_experiment_figures.py::PALETTE:
%     blue    model / flow / learned quantities
%     green   the value term (energy values, L_value)
%     orange  the teacher (universal neural potential)
%     grey    neutral structure
%  Vermilion (= gradient supervision) and purple (= scrambled control)
%  are deliberately ABSENT: neither concept appears in this figure.
%
%  USAGE
%   (a) standalone (verification / slides):   pdflatex fig1_objective.tex
%   (b) inside the paper -- main.tex already \usepackage{tikz} and
%       \usepackage{amsmath}:
%         \input{figures/fig1_objective_float}
%       or, if the writer prefers to own the float:
%         \begin{figure}[t]\centering
%           \input{figures/fig1_objective}
%           \caption{...}\label{fig:objective}
%         \end{figure}
%
%  SIZE.  Natural size 13.90cm x 8.80cm == 5.47in x 3.46in.
%  ICLR is a SINGLE 5.5in text column, so this is a full-text-width
%  figure at scale 1.0.  DO NOT rescale it and DO NOT wrap it in
%  \resizebox: every glyph inside is \small / \footnotesize /
%  \scriptsize of the 10pt body font, and scaling desynchronises
%  in-figure text from the caption.
%
%  LAYOUT.  A 2x2 grid, read left-to-right, top-to-bottom:
%    (a) the generator      composition AND geometry from one flow
%    (b) the teacher        energy values for K displacements of one parent
%    (c) the objective      L_ours = L_FM + lambda_E L_value, and nothing else
%    (d) the scope          two basins of ONE composition; offset free
%  Panel widths 6.80 / 6.78 cm; no wire crosses a panel boundary, so
%  there is no routing to collide with text.
%
%  REVISION C8 (this pass).  Two changes, both confined to panels (c)/(d):
%   * (c) now displays the complete objective, L_ours = L_FM + lambda_E
%     L_value, so the reader cannot infer a hidden force or anchor branch.
%     Neither branch was ever drawn here and neither is added.
%   * (d) previously drew two groups of two DIFFERENT compositions
%     (c_A, c_B).  That was wrong twice over.  It misstated the failure
%     mode -- the unidentified offset is already present WITHIN a single
%     conditional density -- and it flirted with the claim that one
%     composition contributes exactly one group, which the composition
%     audit refutes: in the training perturbation shard
%     (processed_data/omol25_4m_processed/perturbation_train_n30000_s0.pt)
%     30,000 parent groups carry only
%     22,822 distinct (element multiset, total charge) keys, 2,204 keys
%     are shared by more than one parent (9,382 parents, 31%; max
%     multiplicity 80).  Panel (d) therefore now shows two groups A and B
%     of the SAME composition c: two disconnected supervised regions of
%     one density p_theta(.|c), each with its own free offset b_A, b_B.
%     "ordering fixed" / "pinned" wording is replaced throughout by
%     "the local relation between log-density and energy is constrained".
%  In the EVALUATION population (val_data_processed.pt; the 120 grid
%  parents and the 93 primary parents) every parent has a distinct
%  composition, multiplicity 1 -- that is an audited property of the
%  evaluation set, not of the supervision design, and this figure does
%  not assert it.
% =====================================================================

% ---- standalone bootstrap (skipped when \input into a tikz-enabled doc) ----
\ifx\tikzpicture\undefined
  \documentclass[10pt]{article}
  % 0.05cm of slack past the 13.90cm bounding box so that the panel border
  % stroke width does not report a spurious overfull \hbox in standalone mode.
  \usepackage[paperwidth=14.40cm,paperheight=9.20cm,margin=0.20cm,noheadfoot]{geometry}
  \usepackage{amsmath,amssymb,bm}
  \usepackage{times}
  \usepackage{tikz}
  \setlength{\topskip}{0pt}
  \setlength{\parindent}{0pt}
  \pagestyle{empty}
  \begin{document}%
  \def\figOneFinish{\end{document}}
\fi
\providecommand{\figOneFinish}{\relax}
\providecommand{\figOneScale}{1.0}

\usetikzlibrary{positioning,calc,fit,backgrounds,arrows,shapes.geometric}

% ---------------------------------------------------------------- palette
% Okabe & Ito (2008), CVD-safe qualitative set.
\providecommand{\oiOrangeDefined}{}
\definecolor{oiOrange}{HTML}{E69F00}
\definecolor{oiSky}   {HTML}{56B4E9}
\definecolor{oiGreen} {HTML}{009E73}
\definecolor{oiBlue}  {HTML}{0072B2}
\definecolor{oiPurple}{HTML}{CC79A7}
\definecolor{oiGrey}  {HTML}{7F7F7F}
\definecolor{oiGreyL} {HTML}{C4C4C4}
\definecolor{oiGreyD} {HTML}{4D4D4D}

% ---------------------------------------------------------------- type
% Body 10pt Times => \small 9pt, \footnotesize 8pt, \scriptsize 7pt.
% Titles are 9pt bold; all explanatory text is 8pt; 7pt is used only for
% short qualifiers.  Nothing in this figure is smaller than 7pt.
\tikzset{
  fpanel/.style   ={rounded corners=2.4pt, draw=#1!55!white, line width=0.7pt,
                    fill=#1!4},
  ftitle/.style   ={anchor=north west, font=\small\bfseries, inner sep=0pt,
                    text=oiGreyD},
  fsub/.style     ={anchor=north west, font=\footnotesize, inner sep=0pt,
                    text=oiGrey, align=left},
  fbody/.style    ={anchor=north west, font=\footnotesize, inner sep=0pt,
                    align=left},
  fnote/.style    ={anchor=north west, font=\scriptsize, inner sep=0pt,
                    align=left, text=oiGreyD},
  fctr/.style     ={anchor=north, font=\footnotesize, inner sep=0pt,
                    align=center},
  fctrn/.style    ={anchor=north, font=\scriptsize, inner sep=0pt,
                    align=center, text=oiGreyD},
  flow/.style     ={-stealth', line width=1.0pt, draw=oiBlue},
  teach/.style    ={-stealth', line width=1.0pt, draw=oiOrange!88!black},
  freearr/.style  ={stealth'-stealth', line width=0.8pt, draw=oiGrey,
                    dash pattern=on 2.4pt off 1.8pt},
  eqbox/.style    ={rounded corners=2.4pt, draw=oiGreen!65!white,
                    line width=0.8pt, fill=oiGreen!8},
}

% ------------------------------------------------------- molecule glyphs
% Deterministic coordinates: the figure is identical on every compile.
\newcommand{\fAtom}[3]{\fill[#3] (#1,#2) circle (2.05pt);%
  \draw[oiGreyD!55, line width=0.2pt] (#1,#2) circle (2.05pt);}
\newcommand{\fAtomS}[3]{\fill[#3] (#1,#2) circle (1.35pt);%
  \draw[oiGreyD!45, line width=0.15pt] (#1,#2) circle (1.35pt);}
\newcommand{\fBond}[4]{\draw[oiGrey!80, line width=0.7pt] (#1,#2) -- (#3,#4);}

% t = 0 : unstructured cloud, no bonds, no identities
\newcommand{\fMolNoise}{%
  \fAtomS{-0.36}{ 0.32}{oiGreyL}\fAtomS{ 0.12}{ 0.38}{oiGreyL}%
  \fAtomS{ 0.38}{ 0.06}{oiGreyL}\fAtomS{-0.10}{-0.02}{oiGreyL}%
  \fAtomS{-0.42}{-0.28}{oiGreyL}\fAtomS{ 0.06}{-0.36}{oiGreyL}%
  \fAtomS{ 0.42}{-0.34}{oiGreyL}\fAtomS{ 0.26}{ 0.24}{oiGreyL}%
  \fAtomS{-0.18}{ 0.12}{oiGreyL}}

% t = 1 : a structure with distinguishable element identities.  Element
% colours are neutral (grey / blue / sky / purple) so they cannot be read
% as the semantic green (value) or orange (teacher).
\newcommand{\fMolFinal}{%
  \fBond{ 0.000}{ 0.30}{-0.285}{ 0.09}\fBond{-0.285}{ 0.09}{-0.176}{-0.25}%
  \fBond{-0.176}{-0.25}{ 0.176}{-0.25}\fBond{ 0.176}{-0.25}{ 0.285}{ 0.09}%
  \fBond{ 0.285}{ 0.09}{ 0.000}{ 0.30}%
  \fBond{ 0.285}{ 0.09}{ 0.600}{ 0.27}%
  \fBond{-0.285}{ 0.09}{-0.560}{ 0.32}%
  \fBond{-0.176}{-0.25}{-0.385}{-0.47}%
  \fAtom{ 0.000}{ 0.30}{oiGreyD!75}\fAtom{-0.285}{ 0.09}{oiGreyD!75}%
  \fAtom{-0.176}{-0.25}{oiGreyD!75}\fAtom{ 0.176}{-0.25}{oiGreyD!75}%
  \fAtom{ 0.285}{ 0.09}{oiBlue}%
  \fAtom{ 0.600}{ 0.27}{oiPurple}%
  \fAtomS{-0.560}{ 0.32}{oiSky}%
  \fAtomS{-0.385}{-0.47}{white}}

% the same structure, small, for the displacement cloud in panel (b)
\newcommand{\fMolSmall}[1]{%
  \draw[#1, line width=0.55pt]
    ( 0.000, 0.24) -- (-0.228, 0.07) -- (-0.141,-0.20) --
    ( 0.141,-0.20) -- ( 0.228, 0.07) -- cycle;
  \draw[#1, line width=0.55pt] ( 0.228, 0.07) -- ( 0.470, 0.22);
  \draw[#1, line width=0.55pt] (-0.228, 0.07) -- (-0.440, 0.25);
  \foreach \px/\py in {0/0.24, -0.228/0.07, -0.141/-0.20, 0.141/-0.20,
                       0.228/0.07, 0.470/0.22, -0.440/0.25}
    {\fill[#1] (\px,\py) circle (1.15pt);}}

% =====================================================================
\begin{tikzpicture}[x=1cm, y=1cm, scale=\figOneScale, >=stealth']
% stable bounding box
\path (0,0) rectangle (13.90,8.80);

% ---- panel frames ----------------------------------------------------
\draw[fpanel=oiBlue]   (0.02,4.85) rectangle (6.82,8.75);
\draw[fpanel=oiOrange] (7.10,4.85) rectangle (13.88,8.75);
\draw[fpanel=oiGreen]  (0.02,0.05) rectangle (6.82,4.55);
\draw[fpanel=oiGrey]   (7.10,0.05) rectangle (13.88,4.55);

% =====================================================================
% (a) THE GENERATOR
% =====================================================================
\node[ftitle] at (0.22,8.66) {(a) A de novo flow};
\node[fsub, text width=6.30cm] at (0.22,8.30)
     {composition \emph{and} geometry from one model};

% prior
\begin{scope}[shift={(1.15,7.34)}]\fMolNoise\end{scope}
\node[fctr, text width=2.60cm] at (1.32,6.66)
     {$x_0\sim\mathcal{N}(0,\sigma^{2}I)$};
\node[fctrn, text width=2.40cm] at (1.32,6.28) {random types};

% flow arrow
\draw[flow] (2.20,7.34) -- (3.90,7.34);
\node[anchor=south, font=\footnotesize, text=oiBlue!80!black] at (3.05,7.46)
     {$v_\theta(x_t,t)$};
\node[anchor=north, font=\scriptsize, text=oiGreyD] at (3.05,7.20)
     {flow matching};

% sample
\begin{scope}[shift={(5.15,7.34)}]\fMolFinal\end{scope}
\node[fctr, text width=2.60cm] at (5.15,6.66) {$(c,\,x)$};
\node[fctrn, text width=3.10cm] at (5.15,6.30)
     {atom types $c$, coordinates $x$};

\node[fbody, text width=6.46cm] at (0.22,5.94)
     {No bond labels anywhere. With the discrete labels
      clamped, reversing the positional ODE gives
      $\log q_\theta(x\mid c)$.};

% =====================================================================
% (b) THE TEACHER
% =====================================================================
\node[ftitle] at (7.30,8.66) {(b) Teacher energy values};
\node[fsub, text width=6.30cm] at (7.30,8.30)
     {$K$ displaced geometries of \emph{one} parent};

% displacement cloud: two faint copies + one dark reference structure
\draw[oiGrey!45, line width=0.5pt, dash pattern=on 1.8pt off 1.6pt]
      (8.35,7.24) circle (0.72);
\begin{scope}[shift={(8.19,7.40)}]\fMolSmall{oiGreyL}\end{scope}
\begin{scope}[shift={(8.54,7.10)}]\fMolSmall{oiGreyL}\end{scope}
\begin{scope}[shift={(8.33,7.22)}]\fMolSmall{oiGreyD!80}\end{scope}
\node[fctr, text width=2.40cm] at (8.35,6.36) {$x_1,\dots,x_K$};
\node[fctrn, text width=2.60cm] at (8.35,6.04) {same composition $c$};

% teacher arrow
\draw[teach] (9.28,7.24) -- (10.22,7.24);
\node[anchor=south, font=\footnotesize, text=oiOrange!55!black] at (9.75,7.36)
     {$E(c,\cdot)$};

% energy ladder
\draw[-stealth', line width=0.7pt, draw=oiGreyD] (10.72,6.42) -- (10.72,7.94);
\node[anchor=east, font=\footnotesize, text=oiGreyD] at (10.62,7.84) {$E$};
\foreach \yy in {6.58, 6.86, 7.12, 7.44, 7.74}
  {\draw[oiGrey!70, line width=0.5pt] (10.72,\yy) -- (11.06,\yy);
   \fill[oiOrange!85!black] (11.06,\yy) circle (1.6pt);}
\node[anchor=west, font=\footnotesize] at (11.24,7.62)
     {$E_k=E(c,x_k)$};
\node[anchor=north west, font=\scriptsize, text width=2.50cm, text=oiGreyD]
      at (11.24,7.30) {one value per geometry};

\node[fbody, text width=6.34cm] at (7.30,5.62)
     {A universal neural potential over $83$ elements supplies every
      value. Only the values are used.};

% =====================================================================
% (c) THE OBJECTIVE
% =====================================================================
\node[ftitle] at (0.22,4.46) {(c) Proposed energy-value supervision};

\draw[eqbox] (0.16,2.86) rectangle (6.68,4.06);
\node[anchor=center, font=\small] at (3.42,3.86)
     {$\mathcal{L}_{\mathrm{ours}}
       =\mathcal{L}_{\mathrm{FM}}+\lambda_E\,\mathcal{L}_{\mathrm{value}}$};
\node[anchor=center, font=\small] at (3.42,3.48)
     {$\mathcal{L}_{\mathrm{value}}
       =\mathrm{Var}_{k}\big[\,\log q_\theta(x_k\mid c)+E_k/kT\,\big]$};
% The cancellation is the selling point of the objective, so it is drawn in
% the value colour rather than in neutral grey.  It is a statement about the
% objective, not a measured claim.
\node[anchor=north, font=\scriptsize, inner sep=0pt, align=center,
      text=oiGreen!45!black, text width=6.20cm] at (3.42,3.22)
     {one added term; the partition function cancels};

% schematic alignment plot
\draw[oiGreyD, line width=0.7pt] (1.36,2.52) -- (1.36,1.02) -- (3.98,1.02);
\node[rotate=90, anchor=south, font=\footnotesize, text=oiGreyD]
      at (1.20,1.77) {$\log q_\theta$};
\node[anchor=north, font=\footnotesize, text=oiGreyD] at (2.70,0.94)
     {$-E_k/kT$};
% the line the value term asks for: slope one, offset free
\draw[oiGreen!70!black, line width=0.9pt, dash pattern=on 2.6pt off 1.8pt]
      (1.52,1.16) -- (3.82,2.40);
% unconstrained (grey) and aligned (green) geometries
\foreach \px/\py in {1.68/2.12, 2.02/1.24, 2.36/2.32, 2.72/1.42,
                     3.08/1.98, 3.46/1.30, 3.72/2.20}
   {\fill[oiGrey!75] (\px,\py) circle (1.7pt);}
\foreach \px/\py in {1.68/1.22, 2.02/1.42, 2.36/1.50, 2.72/1.76,
                     3.08/1.92, 3.46/2.14, 3.72/2.36}
   {\fill[oiGreen] (\px,\py) circle (1.7pt);}
% legend
\fill[oiGrey!75] (4.42,2.32) circle (1.7pt);
\node[anchor=west, font=\footnotesize, text width=2.00cm] at (4.58,2.32)
     {$\mathcal{L}_{\mathrm{value}}$ large};
\fill[oiGreen] (4.42,1.86) circle (1.7pt);
\node[anchor=west, font=\footnotesize, text width=2.00cm] at (4.58,1.86)
     {$\mathcal{L}_{\mathrm{value}}=0$};
\node[fnote, text width=2.30cm] at (4.30,1.50)
     {one variance per group};

\node[anchor=north, font=\scriptsize, text=oiGreyD] at (3.42,0.56)
     {$\mathcal{L}_{\mathrm{value}}=0$ within a group
      $\iff \log q_\theta=-E_k/kT+b$};

% =====================================================================
% (d) THE SCOPE
% =====================================================================
% Two lines of \small\bfseries; the composition fact that used to sit in a
% separate subtitle is folded into the region label below, which is where the
% single conditional density is named anyway.
\node[ftitle, text width=6.50cm, align=left] at (7.30,4.46)
     {(d) Local supervision and the\\path to global control};

% the enclosing region: ONE conditional density, several supervised groups
\draw[rounded corners=2.8pt, draw=oiGrey!50, line width=0.6pt,
      dash pattern=on 2.2pt off 1.8pt] (7.42,1.02) rectangle (13.56,3.54);
\node[anchor=center, font=\scriptsize, text=oiGreyD, inner sep=1.6pt,
      fill=oiGrey!4] at (10.49,3.54)
     {one density $q_\theta(\cdot\mid c)$ --- same composition $c$};

\foreach \cx/\lbl in {8.86/A, 12.12/B}
{%
  \node[anchor=north, font=\footnotesize] at (\cx,3.36)
       {group $\lbl$ (basin $\lbl$)};
  \draw[oiGreen!60!white, line width=0.8pt, fill=oiGreen!7]
        (\cx,2.38) ellipse (1.28 and 0.56);
  \draw[oiGreen!45!black, line width=0.6pt]
        (\cx-0.88,2.26) -- (\cx-0.44,2.52) -- (\cx,2.36)
        -- (\cx+0.44,2.69) -- (\cx+0.84,2.50);
  \foreach \dx/\yy in {-0.88/2.26, -0.44/2.52, 0.00/2.36,
                        0.44/2.69, 0.84/2.50}
    {\fill[oiGreen] (\cx+\dx,\yy) circle (1.7pt);}
  \node[anchor=north, font=\scriptsize, text=oiGreyD] at (\cx,1.74)
       {local relation constrained};
  \node[anchor=north, font=\footnotesize] at (\cx,1.46)
       {offset $b_{\lbl}$ free};
}

\draw[freearr] (10.16,2.38) -- (10.82,2.38);
\node[anchor=south, font=\footnotesize, text=oiGrey] at (10.49,2.46) {?};

% Same mathematical content as before -- b_B - b_A and the relative mass are
% not identified -- with the constructive half of the statement added: what
% would have to be built to identify them.  No claim is weakened.
% WIDTH BUDGET.  The float sets \figOneScale = 1.000 and the picture is measured
% at that scale (see the measurement note in fig1_objective_float.tex: at 1.000 a
% 600 dpi render has zero ink outside the four panel frames; at 0.86 it has 4085
% such pixels, because tikz's `scale' shrinks COORDINATES but not type).  Panel
% (d) is 6.78cm wide at scale 1.000 and every font size here is absolute.  Two lines of at most ~42 \footnotesize characters is the whole
% budget; anything longer runs out of the panel and into panel (c).  Both qualifications survive the compression: the offset
% DIFFERENCE and the relative mass are the two unidentified quantities, and the
% second line names the extension that would identify them.
\node[anchor=north, font=\footnotesize, align=center]
      at (10.49,0.92)
     {$b_B-b_A$, relative mass: not identified.\\
      Connecting groups would share one offset.};

\end{tikzpicture}

\figOneFinish

%           \caption{...}\label{fig:objective}
%         \end{figure}
%
%  SIZE.  Natural size 13.90cm x 8.80cm == 5.47in x 3.46in.
%  ICLR is a SINGLE 5.5in text column, so this is a full-text-width
%  figure at scale 1.0.  DO NOT rescale it and DO NOT wrap it in
%  \resizebox: every glyph inside is \small / \footnotesize /
%  \scriptsize of the 10pt body font, and scaling desynchronises
%  in-figure text from the caption.
%
%  LAYOUT.  A 2x2 grid, read left-to-right, top-to-bottom:
%    (a) the generator      composition AND geometry from one flow
%    (b) the teacher        energy values for K displacements of one parent
%    (c) the objective      L_ours = L_FM + lambda_E L_value, and nothing else
%    (d) the scope          two basins of ONE composition; offset free
%  Panel widths 6.80 / 6.78 cm; no wire crosses a panel boundary, so
%  there is no routing to collide with text.
%
%  REVISION C8 (this pass).  Two changes, both confined to panels (c)/(d):
%   * (c) now displays the complete objective, L_ours = L_FM + lambda_E
%     L_value, so the reader cannot infer a hidden force or anchor branch.
%     Neither branch was ever drawn here and neither is added.
%   * (d) previously drew two groups of two DIFFERENT compositions
%     (c_A, c_B).  That was wrong twice over.  It misstated the failure
%     mode -- the unidentified offset is already present WITHIN a single
%     conditional density -- and it flirted with the claim that one
%     composition contributes exactly one group, which the composition
%     audit refutes: in the training perturbation shard
%     (processed_data/omol25_4m_processed/perturbation_train_n30000_s0.pt)
%     30,000 parent groups carry only
%     22,822 distinct (element multiset, total charge) keys, 2,204 keys
%     are shared by more than one parent (9,382 parents, 31%; max
%     multiplicity 80).  Panel (d) therefore now shows two groups A and B
%     of the SAME composition c: two disconnected supervised regions of
%     one density p_theta(.|c), each with its own free offset b_A, b_B.
%     "ordering fixed" / "pinned" wording is replaced throughout by
%     "the local relation between log-density and energy is constrained".
%  In the EVALUATION population (val_data_processed.pt; the 120 grid
%  parents and the 93 primary parents) every parent has a distinct
%  composition, multiplicity 1 -- that is an audited property of the
%  evaluation set, not of the supervision design, and this figure does
%  not assert it.
% =====================================================================

% ---- standalone bootstrap (skipped when \input into a tikz-enabled doc) ----
\ifx\tikzpicture\undefined
  \documentclass[10pt]{article}
  % 0.05cm of slack past the 13.90cm bounding box so that the panel border
  % stroke width does not report a spurious overfull \hbox in standalone mode.
  \usepackage[paperwidth=14.40cm,paperheight=9.20cm,margin=0.20cm,noheadfoot]{geometry}
  \usepackage{amsmath,amssymb,bm}
  \usepackage{times}
  \usepackage{tikz}
  \setlength{\topskip}{0pt}
  \setlength{\parindent}{0pt}
  \pagestyle{empty}
  \begin{document}%
  \def\figOneFinish{\end{document}}
\fi
\providecommand{\figOneFinish}{\relax}
\providecommand{\figOneScale}{1.0}

\usetikzlibrary{positioning,calc,fit,backgrounds,arrows,shapes.geometric}

% ---------------------------------------------------------------- palette
% Okabe & Ito (2008), CVD-safe qualitative set.
\providecommand{\oiOrangeDefined}{}
\definecolor{oiOrange}{HTML}{E69F00}
\definecolor{oiSky}   {HTML}{56B4E9}
\definecolor{oiGreen} {HTML}{009E73}
\definecolor{oiBlue}  {HTML}{0072B2}
\definecolor{oiPurple}{HTML}{CC79A7}
\definecolor{oiGrey}  {HTML}{7F7F7F}
\definecolor{oiGreyL} {HTML}{C4C4C4}
\definecolor{oiGreyD} {HTML}{4D4D4D}

% ---------------------------------------------------------------- type
% Body 10pt Times => \small 9pt, \footnotesize 8pt, \scriptsize 7pt.
% Titles are 9pt bold; all explanatory text is 8pt; 7pt is used only for
% short qualifiers.  Nothing in this figure is smaller than 7pt.
\tikzset{
  fpanel/.style   ={rounded corners=2.4pt, draw=#1!55!white, line width=0.7pt,
                    fill=#1!4},
  ftitle/.style   ={anchor=north west, font=\small\bfseries, inner sep=0pt,
                    text=oiGreyD},
  fsub/.style     ={anchor=north west, font=\footnotesize, inner sep=0pt,
                    text=oiGrey, align=left},
  fbody/.style    ={anchor=north west, font=\footnotesize, inner sep=0pt,
                    align=left},
  fnote/.style    ={anchor=north west, font=\scriptsize, inner sep=0pt,
                    align=left, text=oiGreyD},
  fctr/.style     ={anchor=north, font=\footnotesize, inner sep=0pt,
                    align=center},
  fctrn/.style    ={anchor=north, font=\scriptsize, inner sep=0pt,
                    align=center, text=oiGreyD},
  flow/.style     ={-stealth', line width=1.0pt, draw=oiBlue},
  teach/.style    ={-stealth', line width=1.0pt, draw=oiOrange!88!black},
  freearr/.style  ={stealth'-stealth', line width=0.8pt, draw=oiGrey,
                    dash pattern=on 2.4pt off 1.8pt},
  eqbox/.style    ={rounded corners=2.4pt, draw=oiGreen!65!white,
                    line width=0.8pt, fill=oiGreen!8},
}

% ------------------------------------------------------- molecule glyphs
% Deterministic coordinates: the figure is identical on every compile.
\newcommand{\fAtom}[3]{\fill[#3] (#1,#2) circle (2.05pt);%
  \draw[oiGreyD!55, line width=0.2pt] (#1,#2) circle (2.05pt);}
\newcommand{\fAtomS}[3]{\fill[#3] (#1,#2) circle (1.35pt);%
  \draw[oiGreyD!45, line width=0.15pt] (#1,#2) circle (1.35pt);}
\newcommand{\fBond}[4]{\draw[oiGrey!80, line width=0.7pt] (#1,#2) -- (#3,#4);}

% t = 0 : unstructured cloud, no bonds, no identities
\newcommand{\fMolNoise}{%
  \fAtomS{-0.36}{ 0.32}{oiGreyL}\fAtomS{ 0.12}{ 0.38}{oiGreyL}%
  \fAtomS{ 0.38}{ 0.06}{oiGreyL}\fAtomS{-0.10}{-0.02}{oiGreyL}%
  \fAtomS{-0.42}{-0.28}{oiGreyL}\fAtomS{ 0.06}{-0.36}{oiGreyL}%
  \fAtomS{ 0.42}{-0.34}{oiGreyL}\fAtomS{ 0.26}{ 0.24}{oiGreyL}%
  \fAtomS{-0.18}{ 0.12}{oiGreyL}}

% t = 1 : a structure with distinguishable element identities.  Element
% colours are neutral (grey / blue / sky / purple) so they cannot be read
% as the semantic green (value) or orange (teacher).
\newcommand{\fMolFinal}{%
  \fBond{ 0.000}{ 0.30}{-0.285}{ 0.09}\fBond{-0.285}{ 0.09}{-0.176}{-0.25}%
  \fBond{-0.176}{-0.25}{ 0.176}{-0.25}\fBond{ 0.176}{-0.25}{ 0.285}{ 0.09}%
  \fBond{ 0.285}{ 0.09}{ 0.000}{ 0.30}%
  \fBond{ 0.285}{ 0.09}{ 0.600}{ 0.27}%
  \fBond{-0.285}{ 0.09}{-0.560}{ 0.32}%
  \fBond{-0.176}{-0.25}{-0.385}{-0.47}%
  \fAtom{ 0.000}{ 0.30}{oiGreyD!75}\fAtom{-0.285}{ 0.09}{oiGreyD!75}%
  \fAtom{-0.176}{-0.25}{oiGreyD!75}\fAtom{ 0.176}{-0.25}{oiGreyD!75}%
  \fAtom{ 0.285}{ 0.09}{oiBlue}%
  \fAtom{ 0.600}{ 0.27}{oiPurple}%
  \fAtomS{-0.560}{ 0.32}{oiSky}%
  \fAtomS{-0.385}{-0.47}{white}}

% the same structure, small, for the displacement cloud in panel (b)
\newcommand{\fMolSmall}[1]{%
  \draw[#1, line width=0.55pt]
    ( 0.000, 0.24) -- (-0.228, 0.07) -- (-0.141,-0.20) --
    ( 0.141,-0.20) -- ( 0.228, 0.07) -- cycle;
  \draw[#1, line width=0.55pt] ( 0.228, 0.07) -- ( 0.470, 0.22);
  \draw[#1, line width=0.55pt] (-0.228, 0.07) -- (-0.440, 0.25);
  \foreach \px/\py in {0/0.24, -0.228/0.07, -0.141/-0.20, 0.141/-0.20,
                       0.228/0.07, 0.470/0.22, -0.440/0.25}
    {\fill[#1] (\px,\py) circle (1.15pt);}}

% =====================================================================
\begin{tikzpicture}[x=1cm, y=1cm, scale=\figOneScale, >=stealth']
% stable bounding box
\path (0,0) rectangle (13.90,8.80);

% ---- panel frames ----------------------------------------------------
\draw[fpanel=oiBlue]   (0.02,4.85) rectangle (6.82,8.75);
\draw[fpanel=oiOrange] (7.10,4.85) rectangle (13.88,8.75);
\draw[fpanel=oiGreen]  (0.02,0.05) rectangle (6.82,4.55);
\draw[fpanel=oiGrey]   (7.10,0.05) rectangle (13.88,4.55);

% =====================================================================
% (a) THE GENERATOR
% =====================================================================
\node[ftitle] at (0.22,8.66) {(a) A de novo flow};
\node[fsub, text width=6.30cm] at (0.22,8.30)
     {composition \emph{and} geometry from one model};

% prior
\begin{scope}[shift={(1.15,7.34)}]\fMolNoise\end{scope}
\node[fctr, text width=2.60cm] at (1.32,6.66)
     {$x_0\sim\mathcal{N}(0,\sigma^{2}I)$};
\node[fctrn, text width=2.40cm] at (1.32,6.28) {random types};

% flow arrow
\draw[flow] (2.20,7.34) -- (3.90,7.34);
\node[anchor=south, font=\footnotesize, text=oiBlue!80!black] at (3.05,7.46)
     {$v_\theta(x_t,t)$};
\node[anchor=north, font=\scriptsize, text=oiGreyD] at (3.05,7.20)
     {flow matching};

% sample
\begin{scope}[shift={(5.15,7.34)}]\fMolFinal\end{scope}
\node[fctr, text width=2.60cm] at (5.15,6.66) {$(c,\,x)$};
\node[fctrn, text width=3.10cm] at (5.15,6.30)
     {atom types $c$, coordinates $x$};

\node[fbody, text width=6.46cm] at (0.22,5.94)
     {No bond labels anywhere. With the discrete labels
      clamped, reversing the positional ODE gives
      $\log q_\theta(x\mid c)$.};

% =====================================================================
% (b) THE TEACHER
% =====================================================================
\node[ftitle] at (7.30,8.66) {(b) Teacher energy values};
\node[fsub, text width=6.30cm] at (7.30,8.30)
     {$K$ displaced geometries of \emph{one} parent};

% displacement cloud: two faint copies + one dark reference structure
\draw[oiGrey!45, line width=0.5pt, dash pattern=on 1.8pt off 1.6pt]
      (8.35,7.24) circle (0.72);
\begin{scope}[shift={(8.19,7.40)}]\fMolSmall{oiGreyL}\end{scope}
\begin{scope}[shift={(8.54,7.10)}]\fMolSmall{oiGreyL}\end{scope}
\begin{scope}[shift={(8.33,7.22)}]\fMolSmall{oiGreyD!80}\end{scope}
\node[fctr, text width=2.40cm] at (8.35,6.36) {$x_1,\dots,x_K$};
\node[fctrn, text width=2.60cm] at (8.35,6.04) {same composition $c$};

% teacher arrow
\draw[teach] (9.28,7.24) -- (10.22,7.24);
\node[anchor=south, font=\footnotesize, text=oiOrange!55!black] at (9.75,7.36)
     {$E(c,\cdot)$};

% energy ladder
\draw[-stealth', line width=0.7pt, draw=oiGreyD] (10.72,6.42) -- (10.72,7.94);
\node[anchor=east, font=\footnotesize, text=oiGreyD] at (10.62,7.84) {$E$};
\foreach \yy in {6.58, 6.86, 7.12, 7.44, 7.74}
  {\draw[oiGrey!70, line width=0.5pt] (10.72,\yy) -- (11.06,\yy);
   \fill[oiOrange!85!black] (11.06,\yy) circle (1.6pt);}
\node[anchor=west, font=\footnotesize] at (11.24,7.62)
     {$E_k=E(c,x_k)$};
\node[anchor=north west, font=\scriptsize, text width=2.50cm, text=oiGreyD]
      at (11.24,7.30) {one value per geometry};

\node[fbody, text width=6.34cm] at (7.30,5.62)
     {A universal neural potential over $83$ elements supplies every
      value. Only the values are used.};

% =====================================================================
% (c) THE OBJECTIVE
% =====================================================================
\node[ftitle] at (0.22,4.46) {(c) Proposed energy-value supervision};

\draw[eqbox] (0.16,2.86) rectangle (6.68,4.06);
\node[anchor=center, font=\small] at (3.42,3.86)
     {$\mathcal{L}_{\mathrm{ours}}
       =\mathcal{L}_{\mathrm{FM}}+\lambda_E\,\mathcal{L}_{\mathrm{value}}$};
\node[anchor=center, font=\small] at (3.42,3.48)
     {$\mathcal{L}_{\mathrm{value}}
       =\mathrm{Var}_{k}\big[\,\log q_\theta(x_k\mid c)+E_k/kT\,\big]$};
% The cancellation is the selling point of the objective, so it is drawn in
% the value colour rather than in neutral grey.  It is a statement about the
% objective, not a measured claim.
\node[anchor=north, font=\scriptsize, inner sep=0pt, align=center,
      text=oiGreen!45!black, text width=6.20cm] at (3.42,3.22)
     {one added term; the partition function cancels};

% schematic alignment plot
\draw[oiGreyD, line width=0.7pt] (1.36,2.52) -- (1.36,1.02) -- (3.98,1.02);
\node[rotate=90, anchor=south, font=\footnotesize, text=oiGreyD]
      at (1.20,1.77) {$\log q_\theta$};
\node[anchor=north, font=\footnotesize, text=oiGreyD] at (2.70,0.94)
     {$-E_k/kT$};
% the line the value term asks for: slope one, offset free
\draw[oiGreen!70!black, line width=0.9pt, dash pattern=on 2.6pt off 1.8pt]
      (1.52,1.16) -- (3.82,2.40);
% unconstrained (grey) and aligned (green) geometries
\foreach \px/\py in {1.68/2.12, 2.02/1.24, 2.36/2.32, 2.72/1.42,
                     3.08/1.98, 3.46/1.30, 3.72/2.20}
   {\fill[oiGrey!75] (\px,\py) circle (1.7pt);}
\foreach \px/\py in {1.68/1.22, 2.02/1.42, 2.36/1.50, 2.72/1.76,
                     3.08/1.92, 3.46/2.14, 3.72/2.36}
   {\fill[oiGreen] (\px,\py) circle (1.7pt);}
% legend
\fill[oiGrey!75] (4.42,2.32) circle (1.7pt);
\node[anchor=west, font=\footnotesize, text width=2.00cm] at (4.58,2.32)
     {$\mathcal{L}_{\mathrm{value}}$ large};
\fill[oiGreen] (4.42,1.86) circle (1.7pt);
\node[anchor=west, font=\footnotesize, text width=2.00cm] at (4.58,1.86)
     {$\mathcal{L}_{\mathrm{value}}=0$};
\node[fnote, text width=2.30cm] at (4.30,1.50)
     {one variance per group};

\node[anchor=north, font=\scriptsize, text=oiGreyD] at (3.42,0.56)
     {$\mathcal{L}_{\mathrm{value}}=0$ within a group
      $\iff \log q_\theta=-E_k/kT+b$};

% =====================================================================
% (d) THE SCOPE
% =====================================================================
% Two lines of \small\bfseries; the composition fact that used to sit in a
% separate subtitle is folded into the region label below, which is where the
% single conditional density is named anyway.
\node[ftitle, text width=6.50cm, align=left] at (7.30,4.46)
     {(d) Local supervision and the\\path to global control};

% the enclosing region: ONE conditional density, several supervised groups
\draw[rounded corners=2.8pt, draw=oiGrey!50, line width=0.6pt,
      dash pattern=on 2.2pt off 1.8pt] (7.42,1.02) rectangle (13.56,3.54);
\node[anchor=center, font=\scriptsize, text=oiGreyD, inner sep=1.6pt,
      fill=oiGrey!4] at (10.49,3.54)
     {one density $q_\theta(\cdot\mid c)$ --- same composition $c$};

\foreach \cx/\lbl in {8.86/A, 12.12/B}
{%
  \node[anchor=north, font=\footnotesize] at (\cx,3.36)
       {group $\lbl$ (basin $\lbl$)};
  \draw[oiGreen!60!white, line width=0.8pt, fill=oiGreen!7]
        (\cx,2.38) ellipse (1.28 and 0.56);
  \draw[oiGreen!45!black, line width=0.6pt]
        (\cx-0.88,2.26) -- (\cx-0.44,2.52) -- (\cx,2.36)
        -- (\cx+0.44,2.69) -- (\cx+0.84,2.50);
  \foreach \dx/\yy in {-0.88/2.26, -0.44/2.52, 0.00/2.36,
                        0.44/2.69, 0.84/2.50}
    {\fill[oiGreen] (\cx+\dx,\yy) circle (1.7pt);}
  \node[anchor=north, font=\scriptsize, text=oiGreyD] at (\cx,1.74)
       {local relation constrained};
  \node[anchor=north, font=\footnotesize] at (\cx,1.46)
       {offset $b_{\lbl}$ free};
}

\draw[freearr] (10.16,2.38) -- (10.82,2.38);
\node[anchor=south, font=\footnotesize, text=oiGrey] at (10.49,2.46) {?};

% Same mathematical content as before -- b_B - b_A and the relative mass are
% not identified -- with the constructive half of the statement added: what
% would have to be built to identify them.  No claim is weakened.
% WIDTH BUDGET.  The float sets \figOneScale = 1.000 and the picture is measured
% at that scale (see the measurement note in fig1_objective_float.tex: at 1.000 a
% 600 dpi render has zero ink outside the four panel frames; at 0.86 it has 4085
% such pixels, because tikz's `scale' shrinks COORDINATES but not type).  Panel
% (d) is 6.78cm wide at scale 1.000 and every font size here is absolute.  Two lines of at most ~42 \footnotesize characters is the whole
% budget; anything longer runs out of the panel and into panel (c).  Both qualifications survive the compression: the offset
% DIFFERENCE and the relative mass are the two unidentified quantities, and the
% second line names the extension that would identify them.
\node[anchor=north, font=\footnotesize, align=center]
      at (10.49,0.92)
     {$b_B-b_A$, relative mass: not identified.\\
      Connecting groups would share one offset.};

\end{tikzpicture}

\figOneFinish

%           \caption{...}\label{fig:objective}
%         \end{figure}
%
%  SIZE.  Natural size 13.90cm x 8.80cm == 5.47in x 3.46in.
%  ICLR is a SINGLE 5.5in text column, so this is a full-text-width
%  figure at scale 1.0.  DO NOT rescale it and DO NOT wrap it in
%  \resizebox: every glyph inside is \small / \footnotesize /
%  \scriptsize of the 10pt body font, and scaling desynchronises
%  in-figure text from the caption.
%
%  LAYOUT.  A 2x2 grid, read left-to-right, top-to-bottom:
%    (a) the generator      composition AND geometry from one flow
%    (b) the teacher        energy values for K displacements of one parent
%    (c) the objective      L_ours = L_FM + lambda_E L_value, and nothing else
%    (d) the scope          two basins of ONE composition; offset free
%  Panel widths 6.80 / 6.78 cm; no wire crosses a panel boundary, so
%  there is no routing to collide with text.
%
%  REVISION C8 (this pass).  Two changes, both confined to panels (c)/(d):
%   * (c) now displays the complete objective, L_ours = L_FM + lambda_E
%     L_value, so the reader cannot infer a hidden force or anchor branch.
%     Neither branch was ever drawn here and neither is added.
%   * (d) previously drew two groups of two DIFFERENT compositions
%     (c_A, c_B).  That was wrong twice over.  It misstated the failure
%     mode -- the unidentified offset is already present WITHIN a single
%     conditional density -- and it flirted with the claim that one
%     composition contributes exactly one group, which the composition
%     audit refutes: in the training perturbation shard
%     (processed_data/omol25_4m_processed/perturbation_train_n30000_s0.pt)
%     30,000 parent groups carry only
%     22,822 distinct (element multiset, total charge) keys, 2,204 keys
%     are shared by more than one parent (9,382 parents, 31%; max
%     multiplicity 80).  Panel (d) therefore now shows two groups A and B
%     of the SAME composition c: two disconnected supervised regions of
%     one density p_theta(.|c), each with its own free offset b_A, b_B.
%     "ordering fixed" / "pinned" wording is replaced throughout by
%     "the local relation between log-density and energy is constrained".
%  In the EVALUATION population (val_data_processed.pt; the 120 grid
%  parents and the 93 primary parents) every parent has a distinct
%  composition, multiplicity 1 -- that is an audited property of the
%  evaluation set, not of the supervision design, and this figure does
%  not assert it.
% =====================================================================

% ---- standalone bootstrap (skipped when \input into a tikz-enabled doc) ----
\ifx\tikzpicture\undefined
  \documentclass[10pt]{article}
  % 0.05cm of slack past the 13.90cm bounding box so that the panel border
  % stroke width does not report a spurious overfull \hbox in standalone mode.
  \usepackage[paperwidth=14.40cm,paperheight=9.20cm,margin=0.20cm,noheadfoot]{geometry}
  \usepackage{amsmath,amssymb,bm}
  \usepackage{times}
  \usepackage{tikz}
  \setlength{\topskip}{0pt}
  \setlength{\parindent}{0pt}
  \pagestyle{empty}
  \begin{document}%
  \def\figOneFinish{\end{document}}
\fi
\providecommand{\figOneFinish}{\relax}
\providecommand{\figOneScale}{1.0}

\usetikzlibrary{positioning,calc,fit,backgrounds,arrows,shapes.geometric}

% ---------------------------------------------------------------- palette
% Okabe & Ito (2008), CVD-safe qualitative set.
\providecommand{\oiOrangeDefined}{}
\definecolor{oiOrange}{HTML}{E69F00}
\definecolor{oiSky}   {HTML}{56B4E9}
\definecolor{oiGreen} {HTML}{009E73}
\definecolor{oiBlue}  {HTML}{0072B2}
\definecolor{oiPurple}{HTML}{CC79A7}
\definecolor{oiGrey}  {HTML}{7F7F7F}
\definecolor{oiGreyL} {HTML}{C4C4C4}
\definecolor{oiGreyD} {HTML}{4D4D4D}

% ---------------------------------------------------------------- type
% Body 10pt Times => \small 9pt, \footnotesize 8pt, \scriptsize 7pt.
% Titles are 9pt bold; all explanatory text is 8pt; 7pt is used only for
% short qualifiers.  Nothing in this figure is smaller than 7pt.
\tikzset{
  fpanel/.style   ={rounded corners=2.4pt, draw=#1!55!white, line width=0.7pt,
                    fill=#1!4},
  ftitle/.style   ={anchor=north west, font=\small\bfseries, inner sep=0pt,
                    text=oiGreyD},
  fsub/.style     ={anchor=north west, font=\footnotesize, inner sep=0pt,
                    text=oiGrey, align=left},
  fbody/.style    ={anchor=north west, font=\footnotesize, inner sep=0pt,
                    align=left},
  fnote/.style    ={anchor=north west, font=\scriptsize, inner sep=0pt,
                    align=left, text=oiGreyD},
  fctr/.style     ={anchor=north, font=\footnotesize, inner sep=0pt,
                    align=center},
  fctrn/.style    ={anchor=north, font=\scriptsize, inner sep=0pt,
                    align=center, text=oiGreyD},
  flow/.style     ={-stealth', line width=1.0pt, draw=oiBlue},
  teach/.style    ={-stealth', line width=1.0pt, draw=oiOrange!88!black},
  freearr/.style  ={stealth'-stealth', line width=0.8pt, draw=oiGrey,
                    dash pattern=on 2.4pt off 1.8pt},
  eqbox/.style    ={rounded corners=2.4pt, draw=oiGreen!65!white,
                    line width=0.8pt, fill=oiGreen!8},
}

% ------------------------------------------------------- molecule glyphs
% Deterministic coordinates: the figure is identical on every compile.
\newcommand{\fAtom}[3]{\fill[#3] (#1,#2) circle (2.05pt);%
  \draw[oiGreyD!55, line width=0.2pt] (#1,#2) circle (2.05pt);}
\newcommand{\fAtomS}[3]{\fill[#3] (#1,#2) circle (1.35pt);%
  \draw[oiGreyD!45, line width=0.15pt] (#1,#2) circle (1.35pt);}
\newcommand{\fBond}[4]{\draw[oiGrey!80, line width=0.7pt] (#1,#2) -- (#3,#4);}

% t = 0 : unstructured cloud, no bonds, no identities
\newcommand{\fMolNoise}{%
  \fAtomS{-0.36}{ 0.32}{oiGreyL}\fAtomS{ 0.12}{ 0.38}{oiGreyL}%
  \fAtomS{ 0.38}{ 0.06}{oiGreyL}\fAtomS{-0.10}{-0.02}{oiGreyL}%
  \fAtomS{-0.42}{-0.28}{oiGreyL}\fAtomS{ 0.06}{-0.36}{oiGreyL}%
  \fAtomS{ 0.42}{-0.34}{oiGreyL}\fAtomS{ 0.26}{ 0.24}{oiGreyL}%
  \fAtomS{-0.18}{ 0.12}{oiGreyL}}

% t = 1 : a structure with distinguishable element identities.  Element
% colours are neutral (grey / blue / sky / purple) so they cannot be read
% as the semantic green (value) or orange (teacher).
\newcommand{\fMolFinal}{%
  \fBond{ 0.000}{ 0.30}{-0.285}{ 0.09}\fBond{-0.285}{ 0.09}{-0.176}{-0.25}%
  \fBond{-0.176}{-0.25}{ 0.176}{-0.25}\fBond{ 0.176}{-0.25}{ 0.285}{ 0.09}%
  \fBond{ 0.285}{ 0.09}{ 0.000}{ 0.30}%
  \fBond{ 0.285}{ 0.09}{ 0.600}{ 0.27}%
  \fBond{-0.285}{ 0.09}{-0.560}{ 0.32}%
  \fBond{-0.176}{-0.25}{-0.385}{-0.47}%
  \fAtom{ 0.000}{ 0.30}{oiGreyD!75}\fAtom{-0.285}{ 0.09}{oiGreyD!75}%
  \fAtom{-0.176}{-0.25}{oiGreyD!75}\fAtom{ 0.176}{-0.25}{oiGreyD!75}%
  \fAtom{ 0.285}{ 0.09}{oiBlue}%
  \fAtom{ 0.600}{ 0.27}{oiPurple}%
  \fAtomS{-0.560}{ 0.32}{oiSky}%
  \fAtomS{-0.385}{-0.47}{white}}

% the same structure, small, for the displacement cloud in panel (b)
\newcommand{\fMolSmall}[1]{%
  \draw[#1, line width=0.55pt]
    ( 0.000, 0.24) -- (-0.228, 0.07) -- (-0.141,-0.20) --
    ( 0.141,-0.20) -- ( 0.228, 0.07) -- cycle;
  \draw[#1, line width=0.55pt] ( 0.228, 0.07) -- ( 0.470, 0.22);
  \draw[#1, line width=0.55pt] (-0.228, 0.07) -- (-0.440, 0.25);
  \foreach \px/\py in {0/0.24, -0.228/0.07, -0.141/-0.20, 0.141/-0.20,
                       0.228/0.07, 0.470/0.22, -0.440/0.25}
    {\fill[#1] (\px,\py) circle (1.15pt);}}

% =====================================================================
\begin{tikzpicture}[x=1cm, y=1cm, scale=\figOneScale, >=stealth']
% stable bounding box
\path (0,0) rectangle (13.90,8.80);

% ---- panel frames ----------------------------------------------------
\draw[fpanel=oiBlue]   (0.02,4.85) rectangle (6.82,8.75);
\draw[fpanel=oiOrange] (7.10,4.85) rectangle (13.88,8.75);
\draw[fpanel=oiGreen]  (0.02,0.05) rectangle (6.82,4.55);
\draw[fpanel=oiGrey]   (7.10,0.05) rectangle (13.88,4.55);

% =====================================================================
% (a) THE GENERATOR
% =====================================================================
\node[ftitle] at (0.22,8.66) {(a) A de novo flow};
\node[fsub, text width=6.30cm] at (0.22,8.30)
     {composition \emph{and} geometry from one model};

% prior
\begin{scope}[shift={(1.15,7.34)}]\fMolNoise\end{scope}
\node[fctr, text width=2.60cm] at (1.32,6.66)
     {$x_0\sim\mathcal{N}(0,\sigma^{2}I)$};
\node[fctrn, text width=2.40cm] at (1.32,6.28) {random types};

% flow arrow
\draw[flow] (2.20,7.34) -- (3.90,7.34);
\node[anchor=south, font=\footnotesize, text=oiBlue!80!black] at (3.05,7.46)
     {$v_\theta(x_t,t)$};
\node[anchor=north, font=\scriptsize, text=oiGreyD] at (3.05,7.20)
     {flow matching};

% sample
\begin{scope}[shift={(5.15,7.34)}]\fMolFinal\end{scope}
\node[fctr, text width=2.60cm] at (5.15,6.66) {$(c,\,x)$};
\node[fctrn, text width=3.10cm] at (5.15,6.30)
     {atom types $c$, coordinates $x$};

\node[fbody, text width=6.46cm] at (0.22,5.94)
     {No bond labels anywhere. With the discrete labels
      clamped, reversing the positional ODE gives
      $\log q_\theta(x\mid c)$.};

% =====================================================================
% (b) THE TEACHER
% =====================================================================
\node[ftitle] at (7.30,8.66) {(b) Teacher energy values};
\node[fsub, text width=6.30cm] at (7.30,8.30)
     {$K$ displaced geometries of \emph{one} parent};

% displacement cloud: two faint copies + one dark reference structure
\draw[oiGrey!45, line width=0.5pt, dash pattern=on 1.8pt off 1.6pt]
      (8.35,7.24) circle (0.72);
\begin{scope}[shift={(8.19,7.40)}]\fMolSmall{oiGreyL}\end{scope}
\begin{scope}[shift={(8.54,7.10)}]\fMolSmall{oiGreyL}\end{scope}
\begin{scope}[shift={(8.33,7.22)}]\fMolSmall{oiGreyD!80}\end{scope}
\node[fctr, text width=2.40cm] at (8.35,6.36) {$x_1,\dots,x_K$};
\node[fctrn, text width=2.60cm] at (8.35,6.04) {same composition $c$};

% teacher arrow
\draw[teach] (9.28,7.24) -- (10.22,7.24);
\node[anchor=south, font=\footnotesize, text=oiOrange!55!black] at (9.75,7.36)
     {$E(c,\cdot)$};

% energy ladder
\draw[-stealth', line width=0.7pt, draw=oiGreyD] (10.72,6.42) -- (10.72,7.94);
\node[anchor=east, font=\footnotesize, text=oiGreyD] at (10.62,7.84) {$E$};
\foreach \yy in {6.58, 6.86, 7.12, 7.44, 7.74}
  {\draw[oiGrey!70, line width=0.5pt] (10.72,\yy) -- (11.06,\yy);
   \fill[oiOrange!85!black] (11.06,\yy) circle (1.6pt);}
\node[anchor=west, font=\footnotesize] at (11.24,7.62)
     {$E_k=E(c,x_k)$};
\node[anchor=north west, font=\scriptsize, text width=2.50cm, text=oiGreyD]
      at (11.24,7.30) {one value per geometry};

\node[fbody, text width=6.34cm] at (7.30,5.62)
     {A universal neural potential over $83$ elements supplies every
      value. Only the values are used.};

% =====================================================================
% (c) THE OBJECTIVE
% =====================================================================
\node[ftitle] at (0.22,4.46) {(c) Proposed energy-value supervision};

\draw[eqbox] (0.16,2.86) rectangle (6.68,4.06);
\node[anchor=center, font=\small] at (3.42,3.86)
     {$\mathcal{L}_{\mathrm{ours}}
       =\mathcal{L}_{\mathrm{FM}}+\lambda_E\,\mathcal{L}_{\mathrm{value}}$};
\node[anchor=center, font=\small] at (3.42,3.48)
     {$\mathcal{L}_{\mathrm{value}}
       =\mathrm{Var}_{k}\big[\,\log q_\theta(x_k\mid c)+E_k/kT\,\big]$};
% The cancellation is the selling point of the objective, so it is drawn in
% the value colour rather than in neutral grey.  It is a statement about the
% objective, not a measured claim.
\node[anchor=north, font=\scriptsize, inner sep=0pt, align=center,
      text=oiGreen!45!black, text width=6.20cm] at (3.42,3.22)
     {one added term; the partition function cancels};

% schematic alignment plot
\draw[oiGreyD, line width=0.7pt] (1.36,2.52) -- (1.36,1.02) -- (3.98,1.02);
\node[rotate=90, anchor=south, font=\footnotesize, text=oiGreyD]
      at (1.20,1.77) {$\log q_\theta$};
\node[anchor=north, font=\footnotesize, text=oiGreyD] at (2.70,0.94)
     {$-E_k/kT$};
% the line the value term asks for: slope one, offset free
\draw[oiGreen!70!black, line width=0.9pt, dash pattern=on 2.6pt off 1.8pt]
      (1.52,1.16) -- (3.82,2.40);
% unconstrained (grey) and aligned (green) geometries
\foreach \px/\py in {1.68/2.12, 2.02/1.24, 2.36/2.32, 2.72/1.42,
                     3.08/1.98, 3.46/1.30, 3.72/2.20}
   {\fill[oiGrey!75] (\px,\py) circle (1.7pt);}
\foreach \px/\py in {1.68/1.22, 2.02/1.42, 2.36/1.50, 2.72/1.76,
                     3.08/1.92, 3.46/2.14, 3.72/2.36}
   {\fill[oiGreen] (\px,\py) circle (1.7pt);}
% legend
\fill[oiGrey!75] (4.42,2.32) circle (1.7pt);
\node[anchor=west, font=\footnotesize, text width=2.00cm] at (4.58,2.32)
     {$\mathcal{L}_{\mathrm{value}}$ large};
\fill[oiGreen] (4.42,1.86) circle (1.7pt);
\node[anchor=west, font=\footnotesize, text width=2.00cm] at (4.58,1.86)
     {$\mathcal{L}_{\mathrm{value}}=0$};
\node[fnote, text width=2.30cm] at (4.30,1.50)
     {one variance per group};

\node[anchor=north, font=\scriptsize, text=oiGreyD] at (3.42,0.56)
     {$\mathcal{L}_{\mathrm{value}}=0$ within a group
      $\iff \log q_\theta=-E_k/kT+b$};

% =====================================================================
% (d) THE SCOPE
% =====================================================================
% Two lines of \small\bfseries; the composition fact that used to sit in a
% separate subtitle is folded into the region label below, which is where the
% single conditional density is named anyway.
\node[ftitle, text width=6.50cm, align=left] at (7.30,4.46)
     {(d) Local supervision and the\\path to global control};

% the enclosing region: ONE conditional density, several supervised groups
\draw[rounded corners=2.8pt, draw=oiGrey!50, line width=0.6pt,
      dash pattern=on 2.2pt off 1.8pt] (7.42,1.02) rectangle (13.56,3.54);
\node[anchor=center, font=\scriptsize, text=oiGreyD, inner sep=1.6pt,
      fill=oiGrey!4] at (10.49,3.54)
     {one density $q_\theta(\cdot\mid c)$ --- same composition $c$};

\foreach \cx/\lbl in {8.86/A, 12.12/B}
{%
  \node[anchor=north, font=\footnotesize] at (\cx,3.36)
       {group $\lbl$ (basin $\lbl$)};
  \draw[oiGreen!60!white, line width=0.8pt, fill=oiGreen!7]
        (\cx,2.38) ellipse (1.28 and 0.56);
  \draw[oiGreen!45!black, line width=0.6pt]
        (\cx-0.88,2.26) -- (\cx-0.44,2.52) -- (\cx,2.36)
        -- (\cx+0.44,2.69) -- (\cx+0.84,2.50);
  \foreach \dx/\yy in {-0.88/2.26, -0.44/2.52, 0.00/2.36,
                        0.44/2.69, 0.84/2.50}
    {\fill[oiGreen] (\cx+\dx,\yy) circle (1.7pt);}
  \node[anchor=north, font=\scriptsize, text=oiGreyD] at (\cx,1.74)
       {local relation constrained};
  \node[anchor=north, font=\footnotesize] at (\cx,1.46)
       {offset $b_{\lbl}$ free};
}

\draw[freearr] (10.16,2.38) -- (10.82,2.38);
\node[anchor=south, font=\footnotesize, text=oiGrey] at (10.49,2.46) {?};

% Same mathematical content as before -- b_B - b_A and the relative mass are
% not identified -- with the constructive half of the statement added: what
% would have to be built to identify them.  No claim is weakened.
% WIDTH BUDGET.  The float sets \figOneScale = 1.000 and the picture is measured
% at that scale (see the measurement note in fig1_objective_float.tex: at 1.000 a
% 600 dpi render has zero ink outside the four panel frames; at 0.86 it has 4085
% such pixels, because tikz's `scale' shrinks COORDINATES but not type).  Panel
% (d) is 6.78cm wide at scale 1.000 and every font size here is absolute.  Two lines of at most ~42 \footnotesize characters is the whole
% budget; anything longer runs out of the panel and into panel (c).  Both qualifications survive the compression: the offset
% DIFFERENCE and the relative mass are the two unidentified quantities, and the
% second line names the extension that would identify them.
\node[anchor=north, font=\footnotesize, align=center]
      at (10.49,0.92)
     {$b_B-b_A$, relative mass: not identified.\\
      Connecting groups would share one offset.};

\end{tikzpicture}

\figOneFinish

%           \caption{...}\label{fig:objective}
%         \end{figure}
%
%  SIZE.  Natural size 13.90cm x 8.80cm == 5.47in x 3.46in.
%  ICLR is a SINGLE 5.5in text column, so this is a full-text-width
%  figure at scale 1.0.  DO NOT rescale it and DO NOT wrap it in
%  \resizebox: every glyph inside is \small / \footnotesize /
%  \scriptsize of the 10pt body font, and scaling desynchronises
%  in-figure text from the caption.
%
%  LAYOUT.  A 2x2 grid, read left-to-right, top-to-bottom:
%    (a) the generator      composition AND geometry from one flow
%    (b) the teacher        energy values for K displacements of one parent
%    (c) the objective      L_ours = L_FM + lambda_E L_value, and nothing else
%    (d) the scope          two basins of ONE composition; offset free
%  Panel widths 6.80 / 6.78 cm; no wire crosses a panel boundary, so
%  there is no routing to collide with text.
%
%  REVISION C8 (this pass).  Two changes, both confined to panels (c)/(d):
%   * (c) now displays the complete objective, L_ours = L_FM + lambda_E
%     L_value, so the reader cannot infer a hidden force or anchor branch.
%     Neither branch was ever drawn here and neither is added.
%   * (d) previously drew two groups of two DIFFERENT compositions
%     (c_A, c_B).  That was wrong twice over.  It misstated the failure
%     mode -- the unidentified offset is already present WITHIN a single
%     conditional density -- and it flirted with the claim that one
%     composition contributes exactly one group, which the composition
%     audit refutes: in the training perturbation shard
%     (processed_data/omol25_4m_processed/perturbation_train_n30000_s0.pt)
%     30,000 parent groups carry only
%     22,822 distinct (element multiset, total charge) keys, 2,204 keys
%     are shared by more than one parent (9,382 parents, 31%; max
%     multiplicity 80).  Panel (d) therefore now shows two groups A and B
%     of the SAME composition c: two disconnected supervised regions of
%     one density p_theta(.|c), each with its own free offset b_A, b_B.
%     "ordering fixed" / "pinned" wording is replaced throughout by
%     "the local relation between log-density and energy is constrained".
%  In the EVALUATION population (val_data_processed.pt; the 120 grid
%  parents and the 93 primary parents) every parent has a distinct
%  composition, multiplicity 1 -- that is an audited property of the
%  evaluation set, not of the supervision design, and this figure does
%  not assert it.
% =====================================================================

% ---- standalone bootstrap (skipped when \input into a tikz-enabled doc) ----
\ifx\tikzpicture\undefined
  \documentclass[10pt]{article}
  % 0.05cm of slack past the 13.90cm bounding box so that the panel border
  % stroke width does not report a spurious overfull \hbox in standalone mode.
  \usepackage[paperwidth=14.40cm,paperheight=9.20cm,margin=0.20cm,noheadfoot]{geometry}
  \usepackage{amsmath,amssymb,bm}
  \usepackage{times}
  \usepackage{tikz}
  \setlength{\topskip}{0pt}
  \setlength{\parindent}{0pt}
  \pagestyle{empty}
  \begin{document}%
  \def\figOneFinish{\end{document}}
\fi
\providecommand{\figOneFinish}{\relax}
\providecommand{\figOneScale}{1.0}

\usetikzlibrary{positioning,calc,fit,backgrounds,arrows,shapes.geometric}

% ---------------------------------------------------------------- palette
% Okabe & Ito (2008), CVD-safe qualitative set.
\providecommand{\oiOrangeDefined}{}
\definecolor{oiOrange}{HTML}{E69F00}
\definecolor{oiSky}   {HTML}{56B4E9}
\definecolor{oiGreen} {HTML}{009E73}
\definecolor{oiBlue}  {HTML}{0072B2}
\definecolor{oiPurple}{HTML}{CC79A7}
\definecolor{oiGrey}  {HTML}{7F7F7F}
\definecolor{oiGreyL} {HTML}{C4C4C4}
\definecolor{oiGreyD} {HTML}{4D4D4D}

% ---------------------------------------------------------------- type
% Body 10pt Times => \small 9pt, \footnotesize 8pt, \scriptsize 7pt.
% Titles are 9pt bold; all explanatory text is 8pt; 7pt is used only for
% short qualifiers.  Nothing in this figure is smaller than 7pt.
\tikzset{
  fpanel/.style   ={rounded corners=2.4pt, draw=#1!55!white, line width=0.7pt,
                    fill=#1!4},
  ftitle/.style   ={anchor=north west, font=\small\bfseries, inner sep=0pt,
                    text=oiGreyD},
  fsub/.style     ={anchor=north west, font=\footnotesize, inner sep=0pt,
                    text=oiGrey, align=left},
  fbody/.style    ={anchor=north west, font=\footnotesize, inner sep=0pt,
                    align=left},
  fnote/.style    ={anchor=north west, font=\scriptsize, inner sep=0pt,
                    align=left, text=oiGreyD},
  fctr/.style     ={anchor=north, font=\footnotesize, inner sep=0pt,
                    align=center},
  fctrn/.style    ={anchor=north, font=\scriptsize, inner sep=0pt,
                    align=center, text=oiGreyD},
  flow/.style     ={-stealth', line width=1.0pt, draw=oiBlue},
  teach/.style    ={-stealth', line width=1.0pt, draw=oiOrange!88!black},
  freearr/.style  ={stealth'-stealth', line width=0.8pt, draw=oiGrey,
                    dash pattern=on 2.4pt off 1.8pt},
  eqbox/.style    ={rounded corners=2.4pt, draw=oiGreen!65!white,
                    line width=0.8pt, fill=oiGreen!8},
}

% ------------------------------------------------------- molecule glyphs
% Deterministic coordinates: the figure is identical on every compile.
\newcommand{\fAtom}[3]{\fill[#3] (#1,#2) circle (2.05pt);%
  \draw[oiGreyD!55, line width=0.2pt] (#1,#2) circle (2.05pt);}
\newcommand{\fAtomS}[3]{\fill[#3] (#1,#2) circle (1.35pt);%
  \draw[oiGreyD!45, line width=0.15pt] (#1,#2) circle (1.35pt);}
\newcommand{\fBond}[4]{\draw[oiGrey!80, line width=0.7pt] (#1,#2) -- (#3,#4);}

% t = 0 : unstructured cloud, no bonds, no identities
\newcommand{\fMolNoise}{%
  \fAtomS{-0.36}{ 0.32}{oiGreyL}\fAtomS{ 0.12}{ 0.38}{oiGreyL}%
  \fAtomS{ 0.38}{ 0.06}{oiGreyL}\fAtomS{-0.10}{-0.02}{oiGreyL}%
  \fAtomS{-0.42}{-0.28}{oiGreyL}\fAtomS{ 0.06}{-0.36}{oiGreyL}%
  \fAtomS{ 0.42}{-0.34}{oiGreyL}\fAtomS{ 0.26}{ 0.24}{oiGreyL}%
  \fAtomS{-0.18}{ 0.12}{oiGreyL}}

% t = 1 : a structure with distinguishable element identities.  Element
% colours are neutral (grey / blue / sky / purple) so they cannot be read
% as the semantic green (value) or orange (teacher).
\newcommand{\fMolFinal}{%
  \fBond{ 0.000}{ 0.30}{-0.285}{ 0.09}\fBond{-0.285}{ 0.09}{-0.176}{-0.25}%
  \fBond{-0.176}{-0.25}{ 0.176}{-0.25}\fBond{ 0.176}{-0.25}{ 0.285}{ 0.09}%
  \fBond{ 0.285}{ 0.09}{ 0.000}{ 0.30}%
  \fBond{ 0.285}{ 0.09}{ 0.600}{ 0.27}%
  \fBond{-0.285}{ 0.09}{-0.560}{ 0.32}%
  \fBond{-0.176}{-0.25}{-0.385}{-0.47}%
  \fAtom{ 0.000}{ 0.30}{oiGreyD!75}\fAtom{-0.285}{ 0.09}{oiGreyD!75}%
  \fAtom{-0.176}{-0.25}{oiGreyD!75}\fAtom{ 0.176}{-0.25}{oiGreyD!75}%
  \fAtom{ 0.285}{ 0.09}{oiBlue}%
  \fAtom{ 0.600}{ 0.27}{oiPurple}%
  \fAtomS{-0.560}{ 0.32}{oiSky}%
  \fAtomS{-0.385}{-0.47}{white}}

% the same structure, small, for the displacement cloud in panel (b)
\newcommand{\fMolSmall}[1]{%
  \draw[#1, line width=0.55pt]
    ( 0.000, 0.24) -- (-0.228, 0.07) -- (-0.141,-0.20) --
    ( 0.141,-0.20) -- ( 0.228, 0.07) -- cycle;
  \draw[#1, line width=0.55pt] ( 0.228, 0.07) -- ( 0.470, 0.22);
  \draw[#1, line width=0.55pt] (-0.228, 0.07) -- (-0.440, 0.25);
  \foreach \px/\py in {0/0.24, -0.228/0.07, -0.141/-0.20, 0.141/-0.20,
                       0.228/0.07, 0.470/0.22, -0.440/0.25}
    {\fill[#1] (\px,\py) circle (1.15pt);}}

% =====================================================================
\begin{tikzpicture}[x=1cm, y=1cm, scale=\figOneScale, >=stealth']
% stable bounding box
\path (0,0) rectangle (13.90,8.80);

% ---- panel frames ----------------------------------------------------
\draw[fpanel=oiBlue]   (0.02,4.85) rectangle (6.82,8.75);
\draw[fpanel=oiOrange] (7.10,4.85) rectangle (13.88,8.75);
\draw[fpanel=oiGreen]  (0.02,0.05) rectangle (6.82,4.55);
\draw[fpanel=oiGrey]   (7.10,0.05) rectangle (13.88,4.55);

% =====================================================================
% (a) THE GENERATOR
% =====================================================================
\node[ftitle] at (0.22,8.66) {(a) A de novo flow};
\node[fsub, text width=6.30cm] at (0.22,8.30)
     {composition \emph{and} geometry from one model};

% prior
\begin{scope}[shift={(1.15,7.34)}]\fMolNoise\end{scope}
\node[fctr, text width=2.60cm] at (1.32,6.66)
     {$x_0\sim\mathcal{N}(0,\sigma^{2}I)$};
\node[fctrn, text width=2.40cm] at (1.32,6.28) {random types};

% flow arrow
\draw[flow] (2.20,7.34) -- (3.90,7.34);
\node[anchor=south, font=\footnotesize, text=oiBlue!80!black] at (3.05,7.46)
     {$v_\theta(x_t,t)$};
\node[anchor=north, font=\scriptsize, text=oiGreyD] at (3.05,7.20)
     {flow matching};

% sample
\begin{scope}[shift={(5.15,7.34)}]\fMolFinal\end{scope}
\node[fctr, text width=2.60cm] at (5.15,6.66) {$(c,\,x)$};
\node[fctrn, text width=3.10cm] at (5.15,6.30)
     {atom types $c$, coordinates $x$};

\node[fbody, text width=6.46cm] at (0.22,5.94)
     {No bond labels anywhere. With the discrete labels
      clamped, reversing the positional ODE gives
      $\log q_\theta(x\mid c)$.};

% =====================================================================
% (b) THE TEACHER
% =====================================================================
\node[ftitle] at (7.30,8.66) {(b) Teacher energy values};
\node[fsub, text width=6.30cm] at (7.30,8.30)
     {$K$ displaced geometries of \emph{one} parent};

% displacement cloud: two faint copies + one dark reference structure
\draw[oiGrey!45, line width=0.5pt, dash pattern=on 1.8pt off 1.6pt]
      (8.35,7.24) circle (0.72);
\begin{scope}[shift={(8.19,7.40)}]\fMolSmall{oiGreyL}\end{scope}
\begin{scope}[shift={(8.54,7.10)}]\fMolSmall{oiGreyL}\end{scope}
\begin{scope}[shift={(8.33,7.22)}]\fMolSmall{oiGreyD!80}\end{scope}
\node[fctr, text width=2.40cm] at (8.35,6.36) {$x_1,\dots,x_K$};
\node[fctrn, text width=2.60cm] at (8.35,6.04) {same composition $c$};

% teacher arrow
\draw[teach] (9.28,7.24) -- (10.22,7.24);
\node[anchor=south, font=\footnotesize, text=oiOrange!55!black] at (9.75,7.36)
     {$E(c,\cdot)$};

% energy ladder
\draw[-stealth', line width=0.7pt, draw=oiGreyD] (10.72,6.42) -- (10.72,7.94);
\node[anchor=east, font=\footnotesize, text=oiGreyD] at (10.62,7.84) {$E$};
\foreach \yy in {6.58, 6.86, 7.12, 7.44, 7.74}
  {\draw[oiGrey!70, line width=0.5pt] (10.72,\yy) -- (11.06,\yy);
   \fill[oiOrange!85!black] (11.06,\yy) circle (1.6pt);}
\node[anchor=west, font=\footnotesize] at (11.24,7.62)
     {$E_k=E(c,x_k)$};
\node[anchor=north west, font=\scriptsize, text width=2.50cm, text=oiGreyD]
      at (11.24,7.30) {one value per geometry};

\node[fbody, text width=6.34cm] at (7.30,5.62)
     {A universal neural potential over $83$ elements supplies every
      value. Only the values are used.};

% =====================================================================
% (c) THE OBJECTIVE
% =====================================================================
\node[ftitle] at (0.22,4.46) {(c) Proposed energy-value supervision};

\draw[eqbox] (0.16,2.86) rectangle (6.68,4.06);
\node[anchor=center, font=\small] at (3.42,3.86)
     {$\mathcal{L}_{\mathrm{ours}}
       =\mathcal{L}_{\mathrm{FM}}+\lambda_E\,\mathcal{L}_{\mathrm{value}}$};
\node[anchor=center, font=\small] at (3.42,3.48)
     {$\mathcal{L}_{\mathrm{value}}
       =\mathrm{Var}_{k}\big[\,\log q_\theta(x_k\mid c)+E_k/kT\,\big]$};
% The cancellation is the selling point of the objective, so it is drawn in
% the value colour rather than in neutral grey.  It is a statement about the
% objective, not a measured claim.
\node[anchor=north, font=\scriptsize, inner sep=0pt, align=center,
      text=oiGreen!45!black, text width=6.20cm] at (3.42,3.22)
     {one added term; the partition function cancels};

% schematic alignment plot
\draw[oiGreyD, line width=0.7pt] (1.36,2.52) -- (1.36,1.02) -- (3.98,1.02);
\node[rotate=90, anchor=south, font=\footnotesize, text=oiGreyD]
      at (1.20,1.77) {$\log q_\theta$};
\node[anchor=north, font=\footnotesize, text=oiGreyD] at (2.70,0.94)
     {$-E_k/kT$};
% the line the value term asks for: slope one, offset free
\draw[oiGreen!70!black, line width=0.9pt, dash pattern=on 2.6pt off 1.8pt]
      (1.52,1.16) -- (3.82,2.40);
% unconstrained (grey) and aligned (green) geometries
\foreach \px/\py in {1.68/2.12, 2.02/1.24, 2.36/2.32, 2.72/1.42,
                     3.08/1.98, 3.46/1.30, 3.72/2.20}
   {\fill[oiGrey!75] (\px,\py) circle (1.7pt);}
\foreach \px/\py in {1.68/1.22, 2.02/1.42, 2.36/1.50, 2.72/1.76,
                     3.08/1.92, 3.46/2.14, 3.72/2.36}
   {\fill[oiGreen] (\px,\py) circle (1.7pt);}
% legend
\fill[oiGrey!75] (4.42,2.32) circle (1.7pt);
\node[anchor=west, font=\footnotesize, text width=2.00cm] at (4.58,2.32)
     {$\mathcal{L}_{\mathrm{value}}$ large};
\fill[oiGreen] (4.42,1.86) circle (1.7pt);
\node[anchor=west, font=\footnotesize, text width=2.00cm] at (4.58,1.86)
     {$\mathcal{L}_{\mathrm{value}}=0$};
\node[fnote, text width=2.30cm] at (4.30,1.50)
     {one variance per group};

\node[anchor=north, font=\scriptsize, text=oiGreyD] at (3.42,0.56)
     {$\mathcal{L}_{\mathrm{value}}=0$ within a group
      $\iff \log q_\theta=-E_k/kT+b$};

% =====================================================================
% (d) THE SCOPE
% =====================================================================
% Two lines of \small\bfseries; the composition fact that used to sit in a
% separate subtitle is folded into the region label below, which is where the
% single conditional density is named anyway.
\node[ftitle, text width=6.50cm, align=left] at (7.30,4.46)
     {(d) Local supervision and the\\path to global control};

% the enclosing region: ONE conditional density, several supervised groups
\draw[rounded corners=2.8pt, draw=oiGrey!50, line width=0.6pt,
      dash pattern=on 2.2pt off 1.8pt] (7.42,1.02) rectangle (13.56,3.54);
\node[anchor=center, font=\scriptsize, text=oiGreyD, inner sep=1.6pt,
      fill=oiGrey!4] at (10.49,3.54)
     {one density $q_\theta(\cdot\mid c)$ --- same composition $c$};

\foreach \cx/\lbl in {8.86/A, 12.12/B}
{%
  \node[anchor=north, font=\footnotesize] at (\cx,3.36)
       {group $\lbl$ (basin $\lbl$)};
  \draw[oiGreen!60!white, line width=0.8pt, fill=oiGreen!7]
        (\cx,2.38) ellipse (1.28 and 0.56);
  \draw[oiGreen!45!black, line width=0.6pt]
        (\cx-0.88,2.26) -- (\cx-0.44,2.52) -- (\cx,2.36)
        -- (\cx+0.44,2.69) -- (\cx+0.84,2.50);
  \foreach \dx/\yy in {-0.88/2.26, -0.44/2.52, 0.00/2.36,
                        0.44/2.69, 0.84/2.50}
    {\fill[oiGreen] (\cx+\dx,\yy) circle (1.7pt);}
  \node[anchor=north, font=\scriptsize, text=oiGreyD] at (\cx,1.74)
       {local relation constrained};
  \node[anchor=north, font=\footnotesize] at (\cx,1.46)
       {offset $b_{\lbl}$ free};
}

\draw[freearr] (10.16,2.38) -- (10.82,2.38);
\node[anchor=south, font=\footnotesize, text=oiGrey] at (10.49,2.46) {?};

% Same mathematical content as before -- b_B - b_A and the relative mass are
% not identified -- with the constructive half of the statement added: what
% would have to be built to identify them.  No claim is weakened.
% WIDTH BUDGET.  The float sets \figOneScale = 1.000 and the picture is measured
% at that scale (see the measurement note in fig1_objective_float.tex: at 1.000 a
% 600 dpi render has zero ink outside the four panel frames; at 0.86 it has 4085
% such pixels, because tikz's `scale' shrinks COORDINATES but not type).  Panel
% (d) is 6.78cm wide at scale 1.000 and every font size here is absolute.  Two lines of at most ~42 \footnotesize characters is the whole
% budget; anything longer runs out of the panel and into panel (c).  Both qualifications survive the compression: the offset
% DIFFERENCE and the relative mass are the two unidentified quantities, and the
% second line names the extension that would identify them.
\node[anchor=north, font=\footnotesize, align=center]
      at (10.49,0.92)
     {$b_B-b_A$, relative mass: not identified.\\
      Connecting groups would share one offset.};

\end{tikzpicture}

\figOneFinish
